%%%%%%%%%%%%%%%%%%%%%%%%%%%%%%%%%%%%%%%%%%%%%%%%%%%%%%%%%%%%%%%%%%%%%%%%%%%%%%%
%  Cours de L3 - Calcul differentiel - M. JOURDAIN
%
%  PREAMBULE COMMUN DE L'UE (identique dans les sept chapitres et dans le
%  volume complet L3_S5_COURS_calcul_differentiel.tex, qui inclut chaque
%  chapitre tel quel grace au paquet docmute).
%  Tire du preambule des cours de MP ; ajouts : hyperref (liens invisibles),
%  listings (code Python des chapitres 6 et 7), macros de l'UE.
%%%%%%%%%%%%%%%%%%%%%%%%%%%%%%%%%%%%%%%%%%%%%%%%%%%%%%%%%%%%%%%%%%%%%%%%%%%%%%%
\documentclass[a4paper,twoside]{report}

\usepackage[utf8]{inputenc}
\usepackage[T1]{fontenc}
\usepackage[french]{babel}
\usepackage[table]{xcolor}
\usepackage{amsmath,amssymb,amsfonts}
\usepackage{geometry}
\usepackage{fancyhdr}
\usepackage{tikz}
\usepackage{array}
\usepackage{multirow}                            % cellules fusionnees des tableaux bilans
\usepackage{enumitem}
\usepackage{graphicx}
\usepackage{tabularx}

% ---- paquets optionnels utilises dans l'UE ----
\usepackage{stmaryrd}                            % crochets doubles des intervalles d'entiers
\usepackage{needspace}                           % pas de titre «~Exemples :~» isole en bas de page
\usepackage{tikz-cd}                             % diagrammes commutatifs (chapitres 2 et 3)
\usepackage{pgfplots}                            % courbes et surfaces
\pgfplotsset{compat=1.18}
\usepackage{listings}                            % code Python (chapitres 6 et 7)
\graphicspath{{/home/armandj/Documents/A Documents/image/}}  % chemin des images (fixe)

\usetikzlibrary{arrows.meta,decorations.pathreplacing,patterns,calc}

\geometry{hmargin=2.5cm,vmargin=1.5cm}

%%%%%%%%%%%%%%%%%%%%%%%%%%%%%%%% COULEURS %%%%%%%%%%%%%%%%%%%%%%%%%%%%%%%%%%%%%%%
%  Le code couleur propre a chaque chapitre est decrit dans l'en-tete du
%  chapitre (bloc «~CODE COULEUR~»). Couleurs disponibles :
\definecolor{rougecours}{RGB}{192,0,0}
\definecolor{couleurN}{RGB}{0,112,192}
\definecolor{couleurZ}{RGB}{0,150,120}
\definecolor{couleurD}{RGB}{224,150,0}
\definecolor{couleurQ}{RGB}{150,80,180}
\definecolor{codegreen}{rgb}{0,0.6,0}
\definecolor{codegray}{rgb}{0.5,0.5,0.5}
\definecolor{codepurple}{rgb}{0.58,0,0.82}
\definecolor{backcolour}{rgb}{0.98,0.98,0.98}
% Mot defini, en gras rouge
\newcommand{\mot}[1]{\textcolor{rougecours}{\textbf{#1}}}
% Notation ou formule mise en avant en rouge
\newcommand{\rmath}[1]{\textcolor{rougecours}{#1}}

%%%%%%%%%%%%%%%%%%%%%%%%%%%% EN-TETES ET PIEDS DE PAGE %%%%%%%%%%%%%%%%%%%%%%%%%

\newcommand{\lauteur}{Cours de L3 - Calcul différentiel - M. JOURDAIN}

\pagestyle{fancy}
\fancyhf{}
\fancyhead[LE]{\footnotesize\rightmark}
\fancyhead[RE]{\footnotesize\leftmark}
\fancyhead[LO]{\footnotesize\leftmark}
\fancyhead[RO]{\footnotesize\rightmark}
\fancyfoot[LE,LO]{\thepage}
\fancyfoot[RE,RO]{\footnotesize\lauteur}
\renewcommand{\headrulewidth}{0pt}
\renewcommand{\footrulewidth}{0pt}

\fancypagestyle{plain}{%
  \fancyhf{}%
  \fancyfoot[C]{\thepage}%
  \renewcommand{\headrulewidth}{0pt}}

%%%%%%%%%%%%%%%%%%%%%%%%%%%%%%%% ENVIRONNEMENTS %%%%%%%%%%%%%%%%%%%%%%%%%%%%%%%%

% Cadre pour les propositions
\newenvironment{cadre}%
 {\par\vspace{2mm}\noindent\begin{tabular}{|@{\hspace{4mm}}p{0.93\textwidth}@{\hspace{4mm}}|}
  \hline \\[-2mm]}%
 {\\[2mm] \hline \end{tabular}\par\vspace{3mm}}

\newcommand{\definition}{\textbf{Définition :}~}
\newcommand{\propo}[1]{\textbf{\underline{Proposition #1 :}}~}
\newcommand{\theoreme}[1]{\textbf{Théorème #1 :}~}
\newcommand{\methode}{\textbf{\underline{$\heartsuit$ Méthode :}}~}
% Enonce de niveau prerequis, donne sans demonstration ou avec une preuve courte
\newcommand{\rappel}{\textit{(Rappel.)}~}

% Demonstration
\newenvironment{demo}%
 {\par\needspace{4\baselineskip}\vspace{1mm}\noindent$\circ$ \underline{Démonstration :}~\setlength{\parskip}{3pt plus 1pt}}%
 {\hfill$\square$\par\vspace{2mm}}

% Etape d'une demonstration longue : nouvel alinea, intitule en italique
\newcommand{\etape}[1]{\par\smallskip\noindent$\triangleright$~\textit{#1}~}

% Exemples et remarques : blocs indentes
\newenvironment{exemple}%
 {\par\needspace{6\baselineskip}\vspace{1mm}\begin{list}{}{\leftmargin=1cm \rightmargin=0.5cm}\item[]%
  \textit{Exemple :}~}%
 {\end{list}\vspace{1mm}}

\newenvironment{remarque}%
 {\par\needspace{6\baselineskip}\vspace{1mm}\begin{list}{}{\leftmargin=1cm \rightmargin=0.5cm}\item[]%
  \textit{\underline{Remarque :}}~}%
 {\end{list}\vspace{1mm}}

\newenvironment{remarques}%
 {\par\needspace{6\baselineskip}\vspace{1mm}\begin{list}{}{\leftmargin=1cm \rightmargin=0.5cm}\item[]%
  \textit{\underline{Remarques :}}%
  \begin{itemize}[label=$\bullet$,leftmargin=0.45cm,topsep=3pt,itemsep=4pt]}%
 {\end{itemize}\end{list}\vspace{1mm}}

\newenvironment{exemples}%
 {\par\needspace{6\baselineskip}\vspace{1mm}\begin{list}{}{\leftmargin=1cm \rightmargin=0.5cm}\item[]%
  \textit{Exemples :}%
  \begin{itemize}[label=$\bullet$,leftmargin=0.45cm,topsep=3pt,itemsep=4pt]}%
 {\end{itemize}\end{list}\vspace{1mm}}

% Rediger un raisonnement type, encadre a gauche
\newenvironment{redaction}%
 {\par\vspace{2mm}\begin{center}\begin{tabular}{|@{\hspace{3mm}}p{0.85\textwidth}@{\hspace{3mm}}}}%
 {\end{tabular}\end{center}\vspace{2mm}}

%--- Style du code Python (lstlisting) -----------------------------------------
\lstdefinestyle{mystyle}{
    backgroundcolor=\color{backcolour}, commentstyle=\color{codegreen},
    keywordstyle=\color{magenta}, numberstyle=\tiny\color{codegray},
    stringstyle=\color{codepurple}, basicstyle=\ttfamily\small,
    breakatwhitespace=false, breaklines=true, captionpos=b, keepspaces=true,
    numbers=left, numbersep=5pt, showspaces=false, showstringspaces=false,
    showtabs=false, tabsize=4, language=Python}
\lstset{style=mystyle}

%%%%%%%%%%%%%%%%%%%%%%%%%%%%%%%% MACROS DE L'UE %%%%%%%%%%%%%%%%%%%%%%%%%%%%%%%%

% Ensembles de nombres
\newcommand{\N}{\mathbb{N}}
\newcommand{\Z}{\mathbb{Z}}
\newcommand{\Q}{\mathbb{Q}}
\newcommand{\R}{\mathbb{R}}
\newcommand{\C}{\mathbb{C}}
\newcommand{\K}{\mathbb{K}}                       % corps de base : R ou C
\newcommand{\intn}[2]{\llbracket #1\,,#2 \rrbracket}   % intervalle d'entiers (stmaryrd)

% Applications
\newcommand{\appli}[5]{\begin{array}{r|ccl}                % application sur deux lignes
  #1 : & #2 & \longrightarrow & #3 \\ & #4 & \longmapsto & #5 \end{array}}
\newcommand{\id}{\mathrm{id}}                     % identite
\newcommand{\ee}{\mathrm{e}}                      % exponentielle
\newcommand{\ii}{\mathrm{i}}                      % unite imaginaire
\newcommand{\dd}{\mathrm{d}}                      % differentielle, element d'integration
\newcommand{\limn}{\lim\limits_{n \to +\infty}}
\newcommand{\cvu}{\xrightarrow[]{\ \mathrm{CU}\ }}           % convergence uniforme

% Algebre lineaire
\newcommand{\Ker}{\operatorname{Ker}}
\newcommand{\Imf}{\operatorname{Im}}              % image (\Im est la partie imaginaire)
\newcommand{\rg}{\operatorname{rg}}
\newcommand{\Vect}{\operatorname{Vect}}
\newcommand{\tr}{\operatorname{tr}}               % trace
\newcommand{\tra}[1]{{#1}^{T}}                    % transposee
\newcommand{\Mnp}[2]{\mathcal{M}_{#1,#2}(\K)}      % matrices a #1 lignes et #2 colonnes
\newcommand{\Mn}[1]{\mathcal{M}_{#1}(\K)}          % matrices carrees
\newcommand{\GL}{\mathrm{GL}}                     % groupe lineaire
\newcommand{\Sp}{\operatorname{Sp}}               % spectre
\newcommand{\diag}{\operatorname{diag}}
\newcommand{\Bc}{\mathcal{B}}                     % bases

% Espaces normes (chapitre 1)
\newcommand{\norme}[1]{\lVert #1 \rVert}         % norme
\newcommand{\nop}[1]{|\mkern-1.5mu|\mkern-1.5mu|#1|\mkern-1.5mu|\mkern-1.5mu|}   % norme subordonnee (triple barre)
\newcommand{\ps}[2]{\langle #1, #2 \rangle}       % produit scalaire
\newcommand{\Bo}{B}                               % boule ouverte B(a, r)
\newcommand{\Bf}{B_{f}}                           % boule fermee B_f(a, r)
\newcommand{\Sph}{S}                              % sphere S(a, r)
\newcommand{\adh}[1]{\overline{#1}}               % adherence
\newcommand{\inte}[1]{\mathring{#1}}              % interieur
\newcommand{\Fr}{\operatorname{Fr}}               % frontiere
\newcommand{\Bor}{\mathcal{B}}                    % fonctions bornees B(X, F)
\newcommand{\Cc}{\mathcal{C}}                     % fonctions continues, de classe C^k
\newcommand{\Lc}{\mathcal{L}}                     % applications lineaires continues L(E, F)
\newcommand{\dist}{d}                             % distance a une partie
\newcommand{\diam}{\operatorname{diam}}           % diametre

% Calcul differentiel (chapitres 2 a 5)
\newcommand{\pd}[2]{\frac{\partial #1}{\partial #2}}   % derivee partielle
\newcommand{\dpd}[2]{\dfrac{\partial #1}{\partial #2}}  % derivee partielle, en ligne en style display
\newcommand{\Isom}{\mathrm{Isom}}                 % isomorphismes continus d inverse continu
\newcommand{\grad}{\nabla}                        % gradient
\newcommand{\Du}[1]{D_{#1}}                       % derivee selon un vecteur
\newcommand{\Hs}{H}                               % matrice hessienne
\newcommand{\Jac}{J}                              % matrice jacobienne
\newcommand{\Ou}{\Omega}                          % ouvert de definition

% Une formule en ligne n'est jamais coupee par un retour a la ligne
\relpenalty=10000
\binoppenalty=10000

\setlength{\parindent}{1cm}

% Liens du sommaire, des signets et des renvois : invisibles (charge en dernier)
\usepackage[hidelinks,bookmarksnumbered,pdfauthor={M. JOURDAIN},pdftitle={Cours de L3 - Calcul differentiel}]{hyperref}
%%%%%%%%%%%%%%%%%%%%%%%%%%%%%%%% CODE COULEUR %%%%%%%%%%%%%%%%%%%%%%%%%%%%%%%%%%
%  Code couleur du chapitre 3 (constant d'une figure a l'autre) :
%     couleurN (bleu)   : la fonction et les ensembles qu'elle definit (courbes, surfaces,
%                         images, lignes de niveau)
%     couleurD (orange) : l'approximation lineaire (tangentes, droites et plans tangents)
%                         et les iteres de Newton (trait plein : Newton, tirets : jacobienne figee)
%     couleurZ (vert)   : les gradients et vecteurs normaux
%     couleurQ (violet) : les voisinages et domaines de depart (V, W, boules, carres)
%     rougecours        : les points etudies (solutions, points singuliers)
%     gris              : les reperes
%  Source : sommaire seul (annexe 2, sections 3.1 a 3.3 du sommaire).
%%%%%%%%%%%%%%%%%%%%%%%%%%%%%%%%%%%%%%%%%%%%%%%%%%%%%%%%%%%%%%%%%%%%%%%%%%%%%%%
\begin{document}

% Ce fichier constitue le chapitre 3
\setcounter{chapter}{2}

%===============================================================================
\chapter[Inversion locale, fonctions implicites]{Inversion locale, fonctions implicites}
%===============================================================================

\section{Introduction historique}

\begin{quote}
\emph{On connaît les beaux résultats obtenus par M.~E.~Picard dans l'étude des équations différentielles et des équations aux dérivées partielles, grâce à sa méthode des approximations successives. Cette méthode s'applique également avec une grande facilité à la théorie des fonctions implicites.}
\begin{flushright}Édouard Goursat - \emph{Sur la théorie des fonctions implicites}, Bulletin de la SMF, 1903\end{flushright}
\end{quote}

\subsection{Résoudre une équation en la linéarisant}

Résoudre une équation $f(x) = y$ par une formule est en général impossible, même pour un polynôme : au-delà du degré $4$, il n'existe pas de formule générale par radicaux. On procède par approximations successives. Au \textsc{i}\textsuperscript{er} siècle, \textbf{Héron d'Alexandrie} calcule, dans ses \emph{Métriques}, une valeur approchée de $\sqrt{720}$ : il part de $27$, dont le carré $729$ est proche de $720$, prend la moyenne de $27$ et de $\frac{720}{27}$, et obtient $26 + \frac{5}{6}$, dont le carré ne dépasse $720$ que de $\frac{1}{36}$. On peut recommencer à partir de cette valeur.

En 1669, dans le manuscrit \emph{De analysi}, \textbf{Isaac Newton} résout $y^{3} - 2y - 5 = 0$ en partant de $y = 2$ : il pose $y = 2 + p$, obtient $p^{3} + 6p^{2} + 10p - 1 = 0$, néglige les termes en $p^{2}$ et $p^{3}$, d'où $p \approx 0{,}1$, et recommence. \textbf{Joseph Raphson} publie une version itérative du procédé en 1690, et \textbf{Thomas Simpson} l'écrit en 1740 avec le calcul des fluxions, y compris pour des systèmes.

\subsection{Des fonctions implicites aux espaces de Banach}

Une équation $f(x, y) = 0$ définit-elle $y$ comme fonction de $x$ ? Au \textsc{xix}\textsuperscript{e} siècle, \textbf{Cauchy} le démontre pour des fonctions développables en séries entières, puis \textbf{Ulisse Dini}, dans ses leçons de Pise (1877-1878), donne l'énoncé encore enseigné pour les variables réelles. En 1903, \textbf{Édouard Goursat} remarque que la méthode des approximations successives d'\textbf{Émile Picard}, conçue pour les équations différentielles, démontre ce théorème en quelques lignes : c'est la phrase citée en exergue, et c'est, à la complétude près, le théorème du point fixe du chapitre 1. En 1927, \textbf{Hildebrandt} et \textbf{Graves} transposent le théorème dans des espaces abstraits, et, en 1948, \textbf{Leonid Kantorovich} étend la méthode de Newton aux espaces de Banach.

Ce chapitre suit cette histoire : Newton (section 3.2), inversion locale (section 3.3), fonctions implicites (section 3.4), hyperplan tangent (section 3.5).

\vfill
% ---- f(x, y) = (x - 0.18 y^2, y + 0.22 x^2) sur [-1, 1]^2 ; a = (0.5, 0.5), f(a) = (0.455, 0.555) ;
% ---- df(a) = [[1, -0.18], [0.22, 1]] ; carre de demi-cote 0.125 : ecart entre f et f(a) + df(a)(x - a) <= 0.0045 (calcules)
\begin{center}
\begin{tikzpicture}[x=1.2cm,y=1.2cm]
  \begin{scope}
    \draw[couleurQ!60,step=0.25] (-1,-1) grid (1,1);
    \draw[couleurQ,thick] (-1,-1) rectangle (1,1);
    \fill[couleurQ!25] (0.375,0.375) rectangle (0.625,0.625);
    \draw[couleurQ,thick] (0.375,0.375) rectangle (0.625,0.625);
    \fill[rougecours] (0.5,0.5) circle (1.3pt);
    \node[rougecours,font=\small,anchor=north west,fill=white,inner sep=1pt] at (0.64,0.36) {$a$};
  \end{scope}
  \draw[gray,thick,->,>=Stealth] (1.12,0) -- (1.8,0) node[midway,above,font=\small] {$f$};
  \begin{scope}[shift={(2.9,0)}]
\draw[couleurN] plot[smooth] coordinates {(-1.180,-0.780) (-1.146,-0.680) (-1.115,-0.580) (-1.088,-0.480) (-1.065,-0.380) (-1.045,-0.280) (-1.029,-0.180) (-1.016,-0.080) (-1.007,0.020) (-1.002,0.120) (-1.000,0.220) (-1.002,0.320) (-1.007,0.420) (-1.016,0.520) (-1.029,0.620) (-1.045,0.720) (-1.065,0.820) (-1.088,0.920) (-1.115,1.020) (-1.146,1.120) (-1.180,1.220)};
\draw[couleurN] plot[smooth] coordinates {(-1.180,-0.780) (-1.080,-0.822) (-0.980,-0.859) (-0.880,-0.892) (-0.780,-0.921) (-0.680,-0.945) (-0.580,-0.965) (-0.480,-0.980) (-0.380,-0.991) (-0.280,-0.998) (-0.180,-1.000) (-0.080,-0.998) (0.020,-0.991) (0.120,-0.980) (0.220,-0.965) (0.320,-0.945) (0.420,-0.921) (0.520,-0.892) (0.620,-0.859) (0.720,-0.822) (0.820,-0.780)};
\draw[couleurN] plot[smooth] coordinates {(-0.930,-0.876) (-0.896,-0.776) (-0.865,-0.676) (-0.838,-0.576) (-0.815,-0.476) (-0.795,-0.376) (-0.779,-0.276) (-0.766,-0.176) (-0.757,-0.076) (-0.752,0.024) (-0.750,0.124) (-0.752,0.224) (-0.757,0.324) (-0.766,0.424) (-0.779,0.524) (-0.795,0.624) (-0.815,0.724) (-0.838,0.824) (-0.865,0.924) (-0.896,1.024) (-0.930,1.124)};
\draw[couleurN] plot[smooth] coordinates {(-1.101,-0.530) (-1.001,-0.572) (-0.901,-0.609) (-0.801,-0.642) (-0.701,-0.671) (-0.601,-0.695) (-0.501,-0.715) (-0.401,-0.730) (-0.301,-0.741) (-0.201,-0.748) (-0.101,-0.750) (-0.001,-0.748) (0.099,-0.741) (0.199,-0.730) (0.299,-0.715) (0.399,-0.695) (0.499,-0.671) (0.599,-0.642) (0.699,-0.609) (0.799,-0.572) (0.899,-0.530)};
\draw[couleurN] plot[smooth] coordinates {(-0.680,-0.945) (-0.646,-0.845) (-0.615,-0.745) (-0.588,-0.645) (-0.565,-0.545) (-0.545,-0.445) (-0.529,-0.345) (-0.516,-0.245) (-0.507,-0.145) (-0.502,-0.045) (-0.500,0.055) (-0.502,0.155) (-0.507,0.255) (-0.516,0.355) (-0.529,0.455) (-0.545,0.555) (-0.565,0.655) (-0.588,0.755) (-0.615,0.855) (-0.646,0.955) (-0.680,1.055)};
\draw[couleurN] plot[smooth] coordinates {(-1.045,-0.280) (-0.945,-0.322) (-0.845,-0.359) (-0.745,-0.392) (-0.645,-0.421) (-0.545,-0.445) (-0.445,-0.465) (-0.345,-0.480) (-0.245,-0.491) (-0.145,-0.498) (-0.045,-0.500) (0.055,-0.498) (0.155,-0.491) (0.255,-0.480) (0.355,-0.465) (0.455,-0.445) (0.555,-0.421) (0.655,-0.392) (0.755,-0.359) (0.855,-0.322) (0.955,-0.280)};
\draw[couleurN] plot[smooth] coordinates {(-0.430,-0.986) (-0.396,-0.886) (-0.365,-0.786) (-0.338,-0.686) (-0.315,-0.586) (-0.295,-0.486) (-0.279,-0.386) (-0.266,-0.286) (-0.257,-0.186) (-0.252,-0.086) (-0.250,0.014) (-0.252,0.114) (-0.257,0.214) (-0.266,0.314) (-0.279,0.414) (-0.295,0.514) (-0.315,0.614) (-0.338,0.714) (-0.365,0.814) (-0.396,0.914) (-0.430,1.014)};
\draw[couleurN] plot[smooth] coordinates {(-1.011,-0.030) (-0.911,-0.072) (-0.811,-0.109) (-0.711,-0.142) (-0.611,-0.171) (-0.511,-0.195) (-0.411,-0.215) (-0.311,-0.230) (-0.211,-0.241) (-0.111,-0.248) (-0.011,-0.250) (0.089,-0.248) (0.189,-0.241) (0.289,-0.230) (0.389,-0.215) (0.489,-0.195) (0.589,-0.171) (0.689,-0.142) (0.789,-0.109) (0.889,-0.072) (0.989,-0.030)};
\draw[couleurN] plot[smooth] coordinates {(-0.180,-1.000) (-0.146,-0.900) (-0.115,-0.800) (-0.088,-0.700) (-0.065,-0.600) (-0.045,-0.500) (-0.029,-0.400) (-0.016,-0.300) (-0.007,-0.200) (-0.002,-0.100) (0.000,0.000) (-0.002,0.100) (-0.007,0.200) (-0.016,0.300) (-0.029,0.400) (-0.045,0.500) (-0.065,0.600) (-0.088,0.700) (-0.115,0.800) (-0.146,0.900) (-0.180,1.000)};
\draw[couleurN] plot[smooth] coordinates {(-1.000,0.220) (-0.900,0.178) (-0.800,0.141) (-0.700,0.108) (-0.600,0.079) (-0.500,0.055) (-0.400,0.035) (-0.300,0.020) (-0.200,0.009) (-0.100,0.002) (0.000,0.000) (0.100,0.002) (0.200,0.009) (0.300,0.020) (0.400,0.035) (0.500,0.055) (0.600,0.079) (0.700,0.108) (0.800,0.141) (0.900,0.178) (1.000,0.220)};
\draw[couleurN] plot[smooth] coordinates {(0.070,-0.986) (0.104,-0.886) (0.135,-0.786) (0.162,-0.686) (0.185,-0.586) (0.205,-0.486) (0.221,-0.386) (0.234,-0.286) (0.243,-0.186) (0.248,-0.086) (0.250,0.014) (0.248,0.114) (0.243,0.214) (0.234,0.314) (0.221,0.414) (0.205,0.514) (0.185,0.614) (0.162,0.714) (0.135,0.814) (0.104,0.914) (0.070,1.014)};
\draw[couleurN] plot[smooth] coordinates {(-1.011,0.470) (-0.911,0.428) (-0.811,0.391) (-0.711,0.358) (-0.611,0.329) (-0.511,0.305) (-0.411,0.285) (-0.311,0.270) (-0.211,0.259) (-0.111,0.252) (-0.011,0.250) (0.089,0.252) (0.189,0.259) (0.289,0.270) (0.389,0.285) (0.489,0.305) (0.589,0.329) (0.689,0.358) (0.789,0.391) (0.889,0.428) (0.989,0.470)};
\draw[couleurN] plot[smooth] coordinates {(0.320,-0.945) (0.354,-0.845) (0.385,-0.745) (0.412,-0.645) (0.435,-0.545) (0.455,-0.445) (0.471,-0.345) (0.484,-0.245) (0.493,-0.145) (0.498,-0.045) (0.500,0.055) (0.498,0.155) (0.493,0.255) (0.484,0.355) (0.471,0.455) (0.455,0.555) (0.435,0.655) (0.412,0.755) (0.385,0.855) (0.354,0.955) (0.320,1.055)};
\draw[couleurN] plot[smooth] coordinates {(-1.045,0.720) (-0.945,0.678) (-0.845,0.641) (-0.745,0.608) (-0.645,0.579) (-0.545,0.555) (-0.445,0.535) (-0.345,0.520) (-0.245,0.509) (-0.145,0.502) (-0.045,0.500) (0.055,0.502) (0.155,0.509) (0.255,0.520) (0.355,0.535) (0.455,0.555) (0.555,0.579) (0.655,0.608) (0.755,0.641) (0.855,0.678) (0.955,0.720)};
\draw[couleurN] plot[smooth] coordinates {(0.570,-0.876) (0.604,-0.776) (0.635,-0.676) (0.662,-0.576) (0.685,-0.476) (0.705,-0.376) (0.721,-0.276) (0.734,-0.176) (0.743,-0.076) (0.748,0.024) (0.750,0.124) (0.748,0.224) (0.743,0.324) (0.734,0.424) (0.721,0.524) (0.705,0.624) (0.685,0.724) (0.662,0.824) (0.635,0.924) (0.604,1.024) (0.570,1.124)};
\draw[couleurN] plot[smooth] coordinates {(-1.101,0.970) (-1.001,0.928) (-0.901,0.891) (-0.801,0.858) (-0.701,0.829) (-0.601,0.805) (-0.501,0.785) (-0.401,0.770) (-0.301,0.759) (-0.201,0.752) (-0.101,0.750) (-0.001,0.752) (0.099,0.759) (0.199,0.770) (0.299,0.785) (0.399,0.805) (0.499,0.829) (0.599,0.858) (0.699,0.891) (0.799,0.928) (0.899,0.970)};
\draw[couleurN] plot[smooth] coordinates {(0.820,-0.780) (0.854,-0.680) (0.885,-0.580) (0.912,-0.480) (0.935,-0.380) (0.955,-0.280) (0.971,-0.180) (0.984,-0.080) (0.993,0.020) (0.998,0.120) (1.000,0.220) (0.998,0.320) (0.993,0.420) (0.984,0.520) (0.971,0.620) (0.955,0.720) (0.935,0.820) (0.912,0.920) (0.885,1.020) (0.854,1.120) (0.820,1.220)};
\draw[couleurN] plot[smooth] coordinates {(-1.180,1.220) (-1.080,1.178) (-0.980,1.141) (-0.880,1.108) (-0.780,1.079) (-0.680,1.055) (-0.580,1.035) (-0.480,1.020) (-0.380,1.009) (-0.280,1.002) (-0.180,1.000) (-0.080,1.002) (0.020,1.009) (0.120,1.020) (0.220,1.035) (0.320,1.055) (0.420,1.079) (0.520,1.108) (0.620,1.141) (0.720,1.178) (0.820,1.220)};
    \fill[couleurN!25] plot coordinates {(0.350,0.406) (0.375,0.410) (0.400,0.415) (0.425,0.420) (0.450,0.425) (0.475,0.430) (0.500,0.436) (0.525,0.442) (0.550,0.448) (0.575,0.454) (0.600,0.461) (0.596,0.486) (0.592,0.511) (0.589,0.536) (0.584,0.561) (0.580,0.586) (0.575,0.611) (0.571,0.636) (0.565,0.661) (0.560,0.686) (0.555,0.711) (0.530,0.704) (0.505,0.698) (0.480,0.692) (0.455,0.686) (0.430,0.680) (0.405,0.675) (0.380,0.670) (0.355,0.665) (0.330,0.660) (0.305,0.656) (0.310,0.631) (0.315,0.606) (0.321,0.581) (0.325,0.556) (0.330,0.531) (0.334,0.506) (0.339,0.481) (0.342,0.456) (0.346,0.431)} -- cycle;
    \fill[rougecours] (0.455,0.555) circle (1.3pt);
    \draw[gray] (0.455,0.555) circle (0.24);
    \node[rougecours,font=\small,anchor=south,fill=white,inner sep=1pt] at (0.455,1.27) {$f(a)$};
    \draw[rougecours,thin] (0.455,1.27) -- (0.455,0.6);
  \end{scope}
  \draw[gray] (3.355,0.795) -- (4.566,0.834);
  \draw[gray] (3.355,0.315) -- (4.566,-0.834);
  \begin{scope}[shift={(5.4,0)}]
    \begin{scope}
      \clip (0,0) circle (1.18);
\draw[couleurN,thick] plot[smooth] coordinates {(-0.632,-0.894) (-0.658,-0.707) (-0.687,-0.519) (-0.717,-0.332) (-0.750,-0.144) (-0.785,0.043) (-0.822,0.231) (-0.861,0.418) (-0.902,0.606)};
\draw[couleurN,thick] plot[smooth] coordinates {(-0.632,-0.894) (-0.444,-0.862) (-0.257,-0.827) (-0.069,-0.790) (0.118,-0.750) (0.306,-0.707) (0.493,-0.662) (0.681,-0.615) (0.868,-0.564)};
\draw[couleurN,thick] plot[smooth] coordinates {(-0.257,-0.827) (-0.283,-0.640) (-0.312,-0.452) (-0.342,-0.265) (-0.375,-0.077) (-0.410,0.110) (-0.447,0.298) (-0.486,0.485) (-0.527,0.673)};
\draw[couleurN,thick] plot[smooth] coordinates {(-0.687,-0.519) (-0.499,-0.487) (-0.312,-0.452) (-0.124,-0.415) (0.063,-0.375) (0.251,-0.332) (0.438,-0.287) (0.626,-0.240) (0.813,-0.189)};
\draw[couleurN,thick] plot[smooth] coordinates {(0.118,-0.750) (0.092,-0.563) (0.063,-0.375) (0.033,-0.188) (0.000,0.000) (-0.035,0.188) (-0.072,0.375) (-0.111,0.562) (-0.152,0.750)};
\draw[couleurN,thick] plot[smooth] coordinates {(-0.750,-0.144) (-0.562,-0.112) (-0.375,-0.077) (-0.188,-0.040) (0.000,0.000) (0.188,0.043) (0.375,0.088) (0.562,0.135) (0.750,0.186)};
\draw[couleurN,thick] plot[smooth] coordinates {(0.493,-0.662) (0.467,-0.475) (0.438,-0.287) (0.408,-0.100) (0.375,0.088) (0.340,0.275) (0.303,0.463) (0.264,0.650) (0.223,0.838)};
\draw[couleurN,thick] plot[smooth] coordinates {(-0.822,0.231) (-0.634,0.263) (-0.447,0.298) (-0.259,0.335) (-0.072,0.375) (0.116,0.418) (0.303,0.463) (0.491,0.510) (0.678,0.561)};
\draw[couleurN,thick] plot[smooth] coordinates {(0.868,-0.564) (0.842,-0.377) (0.813,-0.189) (0.783,-0.002) (0.750,0.186) (0.715,0.373) (0.678,0.561) (0.639,0.748) (0.598,0.936)};
\draw[couleurN,thick] plot[smooth] coordinates {(-0.902,0.606) (-0.714,0.638) (-0.527,0.673) (-0.339,0.710) (-0.152,0.750) (0.036,0.793) (0.223,0.838) (0.411,0.885) (0.598,0.936)};
      \draw[couleurD,thick,densely dashed] (-0.615,-0.915) -- (0.885,-0.585) -- (0.615,0.915) -- (-0.885,0.585) -- cycle;
    \end{scope}
    \draw[gray] (0,0) circle (1.18);
    \fill[rougecours] (0,0) circle (1.3pt);
    \node[gray,font=\small,anchor=north] at (0,-1.22) {zoom $\times 6$};
  \end{scope}
\end{tikzpicture}
\end{center}

\noindent L'application $f(x, y) = (x - 0{,}18\,y^{2}, y + 0{,}22\,x^{2})$ courbe le quadrillage violet en un quadrillage bleu, mais, vu de près autour de $f(a)$, l'image du petit carré se confond, à $0{,}005$ près, avec le parallélogramme orange en tirets, son image par $\dd f(a)$. Ce parallélogramme n'étant pas aplati, $f$ est une bijection près de $a$ : c'est l'idée du théorème d'inversion locale.
\vfill
\newpage
%===============================================================================
\section[La méthode de Newton]{La méthode de Newton}
%===============================================================================

Dans tout le chapitre, $E$, $F$ et $G$ sont des espaces vectoriels normés réels, et $U$ est un ouvert de $E$. Contrairement au chapitre 2, la complétude va jouer un rôle essentiel : les énoncés principaux supposent que $E$ et $F$ sont des espaces de Banach, ce qui est le cas de $\R^{p}$ et de $\Mn{n}$. Toutes les normes de $\R^{p}$ étant équivalentes, on y choisit celle qui simplifie les calculs, le plus souvent $\norme{\cdot}_{\infty}$ ou $\norme{\cdot}_{2}$.

\subsection{Dans \texorpdfstring{$\R$}{R} : la méthode des tangentes}

Soit $f : I \to \R$ dérivable sur un intervalle ouvert, et $x_{n} \in I$ une valeur approchée d'une solution de $f(x) = 0$. On remplace $f$ par son approximation affine en $x_{n}$, dont le graphe est la tangente, et l'on résout l'équation \emph{affine}
\[ f(x_{n}) + f'(x_{n})(x - x_{n}) = 0. \]
Si $f'(x_{n}) \neq 0$, sa solution est
\[ \rmath{x_{n + 1} = x_{n} - \frac{f(x_{n})}{f'(x_{n})}} : \]
c'est l'abscisse du point où la tangente coupe l'axe. On recommence à partir de $x_{n + 1}$.

% ---- f(x) = x^2 - 2, x0 = 3 : x1 = 11/6 = 1.8333, f(x1) = 1.3611 ; x2 = 1.4621 (calcules)
\begin{center}
\begin{tikzpicture}[x=3.4cm,y=0.62cm]
  \draw[gray,->] (-0.15,0) -- (3.45,0) node[below,font=\small] {$x$};
  \draw[gray,->] (0,-2.4) -- (0,8.2);
  \draw[couleurN,very thick,domain=0:3.18,samples=60] plot (\x,{\x*\x-2});
  \draw[couleurD,thick] (1.65,-1.1) -- (3.2,8.2);
  \draw[couleurD,thick] (1.4621,0) -- (2.15,2.5222);
  \draw[gray,densely dotted] (3,0) -- (3,7);
  \draw[gray,densely dotted] (1.8333,0) -- (1.8333,1.3611);
  \fill[couleurD] (3,7) circle (1.5pt);
  \fill[couleurD] (1.8333,1.3611) circle (1.5pt);
  \fill[couleurD] (3,0) circle (1.5pt) node[below,font=\small] {$x_{0}$};
  \fill[couleurD] (1.8333,0) circle (1.5pt) node[below right=-1pt,font=\small] {$x_{1}$};
  \fill[couleurD] (1.4621,0) circle (1.5pt);
  \draw[couleurD,->,>=Stealth] (1.12,-1.7) node[left,font=\small] {$x_{2}$} -- (1.445,-0.15);
  \fill[rougecours] (1.4142,0) circle (1.5pt);
  \draw[rougecours,->,>=Stealth] (1.2,2.1) node[above,font=\small] {$\sqrt{2}$} -- (1.40,0.18);
  \node[couleurN,font=\small,anchor=east] at (2.55,7.4) {$y = x^{2} - 2$};
  \draw[gray] (0.02,7) -- (-0.02,7) node[left,font=\scriptsize] {$7$};
\end{tikzpicture}
\end{center}

\noindent Pour $f(x) = x^{2} - 2$ et $x_{0} = 3$, la tangente en orange au point d'abscisse $x_{0}$ coupe l'axe en $x_{1} = \dfrac{11}{6} \approx 1{,}833$, puis la tangente suivante en $x_{2} \approx 1{,}462$, déjà proche de $\sqrt{2} \approx 1{,}414$. La courbe étant convexe, chaque tangente reste sous la courbe, et les itérés décroissent vers la racine, en rouge.

\begin{exemple}
Pour l'équation de Newton, $f(y) = y^{3} - 2y - 5$ et $y_{0} = 2$, la formule donne $y_{1} = 2 - \dfrac{-1}{10} = 2{,}1$ : c'est exactement le $p \approx 0{,}1$ de Newton, puisque $10p - 1$ est la partie linéaire de son équation en $p$. Les itérés suivants sont $y_{2} \approx 2{,}094\,568$, puis $y_{3} \approx 2{,}094\,551\,481\,7$, et les erreurs successives valent environ $5{,}4 \times 10^{-3}$, $1{,}7 \times 10^{-5}$, puis $1{,}6 \times 10^{-10}$. À chaque étape, l'exposant de l'erreur double à peu près, donc le nombre de décimales exactes aussi : c'est la convergence quadratique, que la proposition 3 établira.
\end{exemple}

\subsection{La méthode de Héron}

Pour $a > 0$, appliquons la méthode à $f(x) = x^{2} - a$ : $x_{n + 1} = x_{n} - \dfrac{x_{n}^{2} - a}{2x_{n}}$, c'est-à-dire
\[ \rmath{x_{n + 1} = \frac{1}{2}\left(x_{n} + \frac{a}{x_{n}}\right)}. \]
C'est le calcul de Héron : la moyenne de $x_{n}$ et de $\dfrac{a}{x_{n}}$. Si $x_{n} > \sqrt{a}$, alors $\dfrac{a}{x_{n}} < \sqrt{a}$, et la moyenne de deux valeurs qui encadrent $\sqrt{a}$ en est une meilleure approximation.

\begin{cadre}
\propo{1} \emph{(Méthode de Héron.)} Soit $a > 0$ et $x_{0} > 0$. La suite définie par $x_{n + 1} = \dfrac{1}{2}\left(x_{n} + \dfrac{a}{x_{n}}\right)$ est bien définie, et
\begin{enumerate}[label=(\roman*),topsep=3pt,itemsep=2pt]
  \item pour tout $n \in \N$, $\rmath{x_{n + 1} - \sqrt{a} = \dfrac{(x_{n} - \sqrt{a})^{2}}{2x_{n}}}$ ; en particulier $x_{n} \geqslant \sqrt{a}$ pour $n \geqslant 1$ ;
  \item la suite $(x_{n})_{n \geqslant 1}$ est décroissante et converge vers $\sqrt{a}$ ;
  \item pour $n \geqslant 1$, $0 \leqslant x_{n + 1} - \sqrt{a} \leqslant \dfrac{(x_{n} - \sqrt{a})^{2}}{2\sqrt{a}}$.
\end{enumerate}
\end{cadre}

\begin{demo}
Par récurrence, $x_{n} > 0$ pour tout $n$, donc la suite est définie. (i) On calcule $2x_{n}(x_{n + 1} - \sqrt{a}) = x_{n}^{2} + a - 2\sqrt{a}\,x_{n} = (x_{n} - \sqrt{a})^{2}$. Le membre de droite est positif, donc $x_{n + 1} \geqslant \sqrt{a}$. (ii) Pour $n \geqslant 1$, $x_{n + 1} - x_{n} = \dfrac{a - x_{n}^{2}}{2x_{n}} \leqslant 0$. La suite, décroissante et minorée par $\sqrt{a}$, converge vers un réel $\ell \geqslant \sqrt{a} > 0$, qui vérifie $\ell = \dfrac{1}{2}\left(\ell + \dfrac{a}{\ell}\right)$, d'où $\ell^{2} = a$ et $\ell = \sqrt{a}$. (iii) Dans (i), on minore $2x_{n}$ par $2\sqrt{a}$.
\end{demo}

% ---- rectangles d'aire 2 : 1 x 2, 3/2 x 4/3, 17/12 x 24/17 ; carre sqrt(2) = 1.4142 (calcules)
\begin{center}
\begin{tikzpicture}[scale=2.4]
  \draw[gray,->] (-0.1,0) -- (2.3,0);
  \draw[gray,->] (0,-0.1) -- (0,2.25);
  \draw[couleurD!45,thick] (0,0) rectangle (1,2);
  \draw[couleurD!75,thick] (0,0) rectangle (1.5,1.3333);
  \draw[couleurD,thick] (0,0) rectangle (1.4167,1.4118);
  \draw[rougecours,densely dashed] (0,0) rectangle (1.4142,1.4142);
  \node[couleurD!60,font=\small,anchor=south west] at (0.03,1.98) {$1 \times 2$};
  \node[couleurD!85,font=\small,anchor=west] at (1.52,1.33) {$\dfrac{3}{2} \times \dfrac{4}{3}$};
  \node[couleurD,font=\small,anchor=west] at (1.52,0.75) {$\dfrac{17}{12} \times \dfrac{24}{17}$};
  \draw[couleurD,thin] (1.51,0.75) -- (1.43,0.75);
  \node[rougecours,font=\small,anchor=north west] at (1.42,-0.02) {carré de côté $\sqrt{2}$};
  \foreach \t/\l in {1/1,1.5/\frac{3}{2},2/2} \draw[gray] (\t,0.02) -- (\t,-0.02);
  \node[gray,font=\scriptsize,below] at (1,-0.02) {$1$};
  \node[gray,font=\scriptsize,left] at (-0.02,2) {$2$};
  \draw[gray] (0.02,2) -- (-0.02,2);
\end{tikzpicture}
\end{center}

\noindent Géométriquement, Héron cherche le côté d'un carré d'aire $a = 2$. Il part d'un rectangle d'aire $2$, de côtés $x_{0} = 1$ et $\dfrac{a}{x_{0}} = 2$, et le remplace par le rectangle de même aire dont un côté est la moyenne des deux précédents : $\dfrac{3}{2} \times \dfrac{4}{3}$, puis $\dfrac{17}{12} \times \dfrac{24}{17}$. Les rectangles, en orange de plus en plus soutenu, deviennent très vite indiscernables du carré cherché, en pointillés rouges.

\begin{exemples}
  \item Pour $a = 2$ et $x_{0} = 1$, on obtient $x_{1} = \dfrac{3}{2}$, $x_{2} = \dfrac{17}{12}$, $x_{3} = \dfrac{577}{408}$, $x_{4} = \dfrac{665\,857}{470\,832}$. Les erreurs $x_{n} - \sqrt{2}$ valent environ $8{,}6 \times 10^{-2}$, $2{,}5 \times 10^{-3}$, $2{,}1 \times 10^{-6}$ et $1{,}6 \times 10^{-12}$ : l'exposant double à chaque étape, comme le prévoit (iii).
  \item Pour $a = 720$ et $x_{0} = 27$, la formule redonne $x_{1} = \dfrac{161}{6} = 26 + \dfrac{5}{6}$, dont le carré vaut $\dfrac{25\,921}{36} = 720 + \dfrac{1}{36}$, comme chez Héron. L'étape suivante donne $x_{2} \approx 26{,}832\,815\,73$, dont le carré ne dépasse $720$ que d'environ $2{,}7 \times 10^{-7}$.
\end{exemples}

\begin{remarque}
\emph{Piège.} La convergence de la méthode de Newton n'est en général que \emph{locale}. Pour $f(x) = x^{2} - a$, la convexité donne la convergence depuis tout $x_{0} > 0$, mais pour $f = \arctan$, dont la seule racine est $0$, les itérés partis de $x_{0} = 2$ s'éloignent en changeant de signe à chaque étape : la tangente, presque horizontale loin de $0$, envoie l'itéré très loin de l'autre côté. L'erreur la plus répandue est de croire que la méthode converge dès que la fonction s'annule.
\end{remarque}

\subsection{Dans un espace de Banach}

Soit $f : U \to F$ différentiable. Près de $x_{n}$, on remplace l'équation $f(x) = 0$ par l'équation affine
\[ f(x_{n}) + \dd f(x_{n})\cdot(x - x_{n}) = 0. \]
Pour la résoudre, il faut que $\dd f(x_{n})$ soit inversible, et, pour contrôler la solution, que son inverse soit continu. D'où la définition suivante.

\begin{center}
\begin{tabularx}{15cm}{||X}
\noindent \textbf{\underline{Définition :}} On note $\rmath{\Isom(E, F)}$ l'ensemble des applications $u \in \Lc(E, F)$ bijectives dont la réciproque $u^{-1}$ est continue, c'est-à-dire appartient à $\Lc(F, E)$. On l'appelle ensemble des \mot{isomorphismes} de $E$ sur $F$. Pour $E = F$, on retrouve $\Isom(E, E) = \GL(E)$.
\end{tabularx}
\end{center}

\begin{remarques}
  \item En dimension finie, une application linéaire bijective a une réciproque linéaire, donc continue : $\Isom(\R^{p}, \R^{n})$ est l'ensemble des applications linéaires bijectives, non vide si et seulement si $p = n$. Pour $f : U \subset \R^{p} \to \R^{p}$, la condition $\dd f(a) \in \Isom(\R^{p}, \R^{p})$ signifie que la matrice jacobienne $\Jac_{f}(a)$ est inversible, c'est-à-dire que le jacobien est non nul.
  \item Entre espaces de Banach, la réciproque d'une application linéaire continue bijective est automatiquement continue : c'est le théorème de Banach, qui dépasse le cadre de ce cours. On n'en aura pas besoin, puisque la continuité de l'inverse figure dans la définition.
\end{remarques}

\begin{cadre}
\propo{2} Soit $E$ et $F$ deux espaces de Banach. L'ensemble $\Isom(E, F)$ est un ouvert de $\Lc(E, F)$, et l'application $u \mapsto u^{-1}$ est continue de $\Isom(E, F)$ dans $\Lc(F, E)$.
\end{cadre}

\begin{demo}
Soit $u_{0} \in \Isom(E, F)$. Pour $u \in \Lc(E, F)$, on écrit $u = u_{0} \circ (u_{0}^{-1} \circ u)$. L'application $\Phi : u \mapsto u_{0}^{-1} \circ u$ est linéaire de $\Lc(E, F)$ dans $\Lc(E)$, et continue, puisque $\nop{u_{0}^{-1} \circ u} \leqslant \nop{u_{0}^{-1}}\,\nop{u}$. Si $\Phi(u) \in \GL(E)$, alors $u$ est bijective, d'inverse $\Phi(u)^{-1} \circ u_{0}^{-1}$ continu : $u \in \Isom(E, F)$. Or $\GL(E)$ est un ouvert de $\Lc(E)$ contenant $\Phi(u_{0}) = \id$ (chapitre 1, proposition 16), donc $\Phi^{-1}(\GL(E))$ est un ouvert contenant $u_{0}$ et contenu dans $\Isom(E, F)$. Sur cet ouvert, $u^{-1} = \Phi(u)^{-1} \circ u_{0}^{-1}$ est la composée de $\Phi$, de l'inversion de $\GL(E)$, continue (chapitre 1, proposition 16), et de la composition à droite par $u_{0}^{-1}$, linéaire continue : elle est continue.
\end{demo}

\noindent Si $\dd f(x_{n}) \in \Isom(E, F)$, l'équation affine a pour unique solution
\[ \rmath{x_{n + 1} = x_{n} - \dd f(x_{n})^{-1}\cdot f(x_{n})}. \]
Dans $\R^{p}$, on ne calcule pas l'inverse de la jacobienne : on résout le système linéaire $\Jac_{f}(x_{n})\,\delta = -f(x_{n})$, puis on pose $x_{n + 1} = x_{n} + \delta$.

\begin{cadre}
\propo{3} \emph{(Convergence quadratique.)} Soit $E$ et $F$ deux espaces de Banach, $f : U \to F$ différentiable, et $c \in U$ tel que $f(c) = 0$ et $\dd f(c) \in \Isom(E, F)$. On suppose que $\dd f$ est $K$-lipschitzienne sur une boule $\Bo(c, \rho) \subset U$ : $\nop{\dd f(x) - \dd f(x')} \leqslant K\norme{x - x'}$. Il existe alors $\delta > 0$ et $C \geqslant 0$ tels que, pour tout $x_{0} \in \Bo(c, \delta)$, la suite de Newton soit définie, reste dans $\Bo(c, \delta)$, converge vers $c$, et vérifie
\[ \rmath{\norme{x_{n + 1} - c} \leqslant C\norme{x_{n} - c}^{2}}. \]
\end{cadre}

\begin{demo}
\etape{Contrôle de l'inverse.} L'application $\dd f$ est continue en $c$, et l'inversion est continue sur l'ouvert $\Isom(E, F)$ (proposition 2). Avec $M = 2\nop{\dd f(c)^{-1}}$, il existe donc $\rho' \in \left]0, \rho\right]$ tel que, pour tout $x \in \Bo(c, \rho')$, $\dd f(x) \in \Isom(E, F)$ et $\nop{\dd f(x)^{-1}} \leqslant M$.

\etape{Une étape.} Soit $x_{n} \in \Bo(c, \rho')$. Comme $f(c) = 0$,
\[ x_{n + 1} - c = \dd f(x_{n})^{-1}\cdot\big(f(c) - f(x_{n}) - \dd f(x_{n})\cdot(c - x_{n})\big). \]
La boule est convexe, et la forme raffinée des accroissements finis (chapitre 2, proposition 13), appliquée entre $x_{n}$ et $c$ avec le point de base $x_{n}$, donne
\[ \norme{f(c) - f(x_{n}) - \dd f(x_{n})\cdot(c - x_{n})} \leqslant \sup_{z \in [x_{n}, c]} \nop{\dd f(z) - \dd f(x_{n})}\,\norme{c - x_{n}} \leqslant K\norme{c - x_{n}}^{2}. \]
Ainsi $\norme{x_{n + 1} - c} \leqslant MK\norme{x_{n} - c}^{2}$.

\etape{Conclusion.} On pose $C = MK$ et $\delta = \min\left(\rho', \dfrac{1}{2C}\right)$, avec $\delta = \rho'$ si $C = 0$. Si $x_{n} \in \Bo(c, \delta)$, alors
\[ \norme{x_{n + 1} - c} \leqslant C\delta\,\norme{x_{n} - c} \leqslant \frac{1}{2}\norme{x_{n} - c}, \]
et, par récurrence, la suite reste dans $\Bo(c, \delta)$, et $\norme{x_{n} - c} \leqslant 2^{-n}\norme{x_{0} - c} \to 0$.
\end{demo}

\begin{remarques}
  \item La majoration de l'erreur par $C$ fois le carré de l'erreur précédente explique le doublement des décimales : si $C\norme{x_{n} - c} \leqslant 10^{-k}$, alors $C\norme{x_{n + 1} - c} \leqslant 10^{-2k}$.
  \item L'hypothèse «~$\dd f$ lipschitzienne~» est satisfaite sur toute boule fermée de $\R^{p}$ incluse dans $U$ dès que $f$ est de classe $\Cc^{2}$, notion que le chapitre 4 définira : la différentielle de $\dd f$ y est bornée, et les accroissements finis s'appliquent.
\end{remarques}

\begin{exemple}
Cherchons l'intersection, dans le quart de plan $x > y > 0$, du cercle $x^{2} + y^{2} = 4$ et de l'hyperbole $xy = 1$. Avec $f(x, y) = (x^{2} + y^{2} - 4, xy - 1)$, on a
\[ \Jac_{f}(x, y) = \begin{pmatrix} 2x & 2y \\ y & x \end{pmatrix}, \qquad \det \Jac_{f}(x, y) = 2(x^{2} - y^{2}). \]
La solution exacte est $c = \left(\sqrt{2 + \sqrt{3}}, \sqrt{2 - \sqrt{3}}\right) \approx (1{,}931\,85 ; 0{,}517\,64)$, puisque $x^{2}$ est une racine de $X^{2} - 4X + 1$. En partant de $x_{0} = (2{,}4 ; 1{,}2)$, la méthode donne $x_{1} \approx (2{,}033 ; 0{,}600)$, $x_{2} \approx (1{,}935\,1 ; 0{,}520\,8)$, et les erreurs $\norme{x_{n} - c}_{\infty}$ valent environ $0{,}68$, $0{,}10$, $3{,}3 \times 10^{-3}$, $4{,}2 \times 10^{-6}$, puis $7{,}2 \times 10^{-12}$.
\end{exemple}

% ---- cercle x^2+y^2=4, hyperbole xy=1 ; x0 = (2.4, 1.2) ; droites linearisees en x0 :
% ---- 4.8x + 2.4y = 11.2 et 1.2x + 2.4y = 3.88 ; x1 = (2.0333, 0.6), x2 = (1.9351, 0.5208) ; c = (1.9319, 0.5176)
% ---- erreurs Newton / jacobienne figee (norme infinie), calculees en Python
\begin{center}
\begin{tikzpicture}
  \begin{scope}[scale=2.75,xshift=-0.85cm]
    \begin{scope}
    \clip (0.85,-0.15) rectangle (2.95,1.65);
    \draw[gray,->] (0.9,0) -- (2.92,0) node[below left,font=\small] {$x$};
    \draw[couleurN,very thick] (0,0) circle (2);
    \draw[couleurN,very thick,domain=0.62:2.92,samples=60] plot (\x,{1/\x});
    \draw[couleurD,thick] (2.383,-0.100) -- (1.533,1.600);
    \draw[couleurD,thick] (0.900,1.167) -- (2.900,0.167);
    \end{scope}
    \fill[couleurD] (2.4,1.2) circle (0.9pt) node[above right=-1pt,font=\small] {$x_{0}$};
    \fill[couleurD] (2.0333,0.6) circle (0.9pt) node[above right=-1pt,font=\small] {$x_{1}$};
    \fill[couleurD] (1.9351,0.5208) circle (0.7pt);
    \fill[rougecours] (1.9319,0.5176) circle (0.9pt);
    \draw[rougecours,->,>=Stealth] (1.55,0.25) node[left,font=\small] {$c$} -- (1.91,0.505);
    \node[couleurN,font=\small,anchor=south west] at (1.13,1.66) {$\mathcal{C}_{1}$};
    \node[couleurN,font=\small,anchor=north east] at (1.08,0.80) {$\mathcal{C}_{2}$};
  \end{scope}
  \begin{scope}[xshift=6.95cm,yshift=-0.35cm]
    \begin{axis}[width=8cm,height=5.9cm,ymode=log,xmin=0,xmax=15.5,ymin=1e-14,ymax=3,
      xlabel={$n$},ylabel={$\norme{x_{n} - c}_{\infty}$},label style={font=\small},tick label style={font=\scriptsize},
      ytick={1e-12,1e-9,1e-6,1e-3,1},axis lines=left,clip=false]
      \addplot[couleurD,thick,mark=*,mark size=1.3pt] coordinates {(0,6.824e-01) (1,1.015e-01) (2,3.272e-03) (3,4.195e-06) (4,7.175e-12)};
      \addplot[couleurD,thick,densely dashed,mark=o,mark size=1.3pt] coordinates {(0,6.824e-01) (1,1.015e-01) (2,2.881e-02) (3,8.751e-03) (4,2.732e-03) (5,8.633e-04) (6,2.743e-04) (7,8.738e-05) (8,2.788e-05) (9,8.901e-06) (10,2.843e-06) (11,9.084e-07) (12,2.903e-07) (13,9.276e-08) (14,2.964e-08) (15,9.473e-09)};
      \node[couleurD,font=\small,anchor=west] at (axis cs:4.3,1e-11) {Newton};
      \node[couleurD,font=\small,anchor=south west,align=left] at (axis cs:8.5,3e-5) {jacobienne\\ figée};
    \end{axis}
  \end{scope}
\end{tikzpicture}
\end{center}

\noindent À gauche, le cercle $\mathcal{C}_{1}$ et l'hyperbole $\mathcal{C}_{2}$, en bleu, sont remplacées près de $x_{0}$ par les droites orange d'équations $f_{i}(x_{0}) + \dd f_{i}(x_{0})\cdot(x - x_{0}) = 0$, parallèles aux tangentes aux lignes de niveau de $f_{1}$ et de $f_{2}$ passant par $x_{0}$ ; leur intersection est $x_{1}$, et l'itéré suivant est déjà presque confondu avec la solution $c$, en rouge. À droite, en échelle logarithmique, l'erreur de la méthode de Newton s'effondre, tandis que celle de la méthode à jacobienne figée, présentée ci-dessous, ne diminue que d'un facteur à peu près constant, voisin de $0{,}3$, à chaque étape.

\subsection{La méthode à jacobienne figée}

Calculer $\dd f(x_{n})$ et résoudre un système linéaire à chaque étape coûte cher. La \mot{méthode de Newton à jacobienne figée} garde la différentielle du point de départ $a$ :
\[ \rmath{x_{n + 1} = x_{n} - \dd f(a)^{-1}\cdot\big(f(x_{n}) - y\big)}, \]
écrite ici pour l'équation $f(x) = y$. Elle converge moins vite, mais sa convergence se démontre par le théorème du point fixe, avec des hypothèses qui ne portent que sur $\dd f$ au voisinage de $a$. C'est l'approche algorithmique de l'inversion locale : cette suite construit la solution, et sa convergence prouve que la solution existe. On commence par le cas où $f$ est une perturbation de l'identité.

\begin{cadre}
\propo{4} \emph{(Perturbation de l'identité.)} Soit $E$ un espace de Banach, $r > 0$, $k \in \left[0, 1\right[$, et $\varphi : \Bf(0, r) \to E$ une application $k$-lipschitzienne telle que $\varphi(0) = 0$. On pose $\psi(h) = h + \varphi(h)$.
\begin{enumerate}[label=(\roman*),topsep=3pt,itemsep=2pt]
  \item Pour tous $h$, $h' \in \Bf(0, r)$, $\rmath{\norme{\psi(h) - \psi(h')} \geqslant (1 - k)\norme{h - h'}}$. En particulier, $\psi$ est injective.
  \item Pour tout $z \in E$ tel que $\norme{z} \leqslant (1 - k)r$, l'équation $\psi(h) = z$ a une unique solution $h \in \Bf(0, r)$ ; elle vérifie $\norme{h} \leqslant \dfrac{\norme{z}}{1 - k}$, et elle est la limite de la suite définie par $h_{0} = 0$ et $\rmath{h_{n + 1} = z - \varphi(h_{n})}$, avec $\norme{h_{n} - h} \leqslant k^{n}r$.
\end{enumerate}
\end{cadre}

\begin{demo}
(i) Comme $h - h' = \psi(h) - \psi(h') - (\varphi(h) - \varphi(h'))$, l'inégalité triangulaire donne
\[ \norme{h - h'} \leqslant \norme{\psi(h) - \psi(h')} + k\norme{h - h'}. \]

(ii) On pose $T(h) = z - \varphi(h)$ ; ses points fixes sont les solutions de $\psi(h) = z$. Pour $h \in \Bf(0, r)$, $\norme{T(h)} \leqslant \norme{z} + \norme{\varphi(h) - \varphi(0)} \leqslant (1 - k)r + kr = r$ : $T$ envoie la boule fermée $\Bf(0, r)$ dans elle-même, et elle est $k$-contractante. La boule fermée est une partie fermée non vide de l'espace de Banach $E$, et le théorème du point fixe (chapitre 1, théorème 13) donne un unique point fixe $h$, limite de la suite $(h_{n})$. Comme $\norme{h_{n + 1} - h} = \norme{T(h_{n}) - T(h)} \leqslant k\norme{h_{n} - h}$, on obtient $\norme{h_{n} - h} \leqslant k^{n}\norme{h_{0} - h} \leqslant k^{n}r$. Enfin, (i) avec $h' = 0$ donne $(1 - k)\norme{h} \leqslant \norme{\psi(h)} = \norme{z}$.
\end{demo}

% ---- psi(x, y) = (x + 0.4 sin y, y + 0.4 sin x) ; phi est 0.4-lipschitzienne pour la norme euclidienne
% ---- image du cercle unite : rayon compris entre 0.6325 et 1.3675 (calcule), donc contient le disque de rayon 0.6
\begin{center}
\begin{tikzpicture}[scale=1.9]
  \begin{scope}
    \foreach \r in {0.25,0.5,0.75} \draw[couleurQ] (0,0) circle (\r);
    \draw[couleurQ,thick] (0,0) circle (1);
    \foreach \a in {0,30,...,330} \draw[couleurQ] (0,0) -- (\a:1);
    \node[couleurQ,font=\small,anchor=south west] at (40:1.02) {$\Bf(0, 1)$};
    \fill[rougecours] (0,0) circle (0.9pt);
  \end{scope}
  \draw[gray,->,>=Stealth] (1.45,0.15) to[bend left=20] node[above,font=\small] {$\psi$} (2.35,0.15);
  \begin{scope}[xshift=3.9cm]
\draw[couleurN] plot[smooth cycle] coordinates {(0.250,0.099) (0.258,0.120) (0.264,0.141) (0.267,0.160) (0.269,0.179) (0.269,0.196) (0.266,0.211) (0.262,0.225) (0.256,0.237) (0.247,0.247) (0.237,0.256) (0.225,0.262) (0.211,0.266) (0.196,0.269) (0.179,0.269) (0.160,0.267) (0.141,0.264) (0.120,0.258) (0.099,0.250) (0.077,0.240) (0.054,0.229) (0.031,0.216) (0.008,0.201) (-0.016,0.184) (-0.039,0.167) (-0.062,0.148) (-0.085,0.128) (-0.106,0.106) (-0.128,0.085) (-0.148,0.062) (-0.167,0.039) (-0.184,0.016) (-0.201,-0.008) (-0.216,-0.031) (-0.229,-0.054) (-0.240,-0.077) (-0.250,-0.099) (-0.258,-0.120) (-0.264,-0.141) (-0.267,-0.160) (-0.269,-0.179) (-0.269,-0.196) (-0.266,-0.211) (-0.262,-0.225) (-0.256,-0.237) (-0.247,-0.247) (-0.237,-0.256) (-0.225,-0.262) (-0.211,-0.266) (-0.196,-0.269) (-0.179,-0.269) (-0.160,-0.267) (-0.141,-0.264) (-0.120,-0.258) (-0.099,-0.250) (-0.077,-0.240) (-0.054,-0.229) (-0.031,-0.216) (-0.008,-0.201) (0.016,-0.184) (0.039,-0.167) (0.062,-0.148) (0.085,-0.128) (0.106,-0.106) (0.128,-0.085) (0.148,-0.062) (0.167,-0.039) (0.184,-0.016) (0.201,0.008) (0.216,0.031) (0.229,0.054) (0.240,0.077)};
\draw[couleurN] plot[smooth cycle] coordinates {(0.500,0.192) (0.516,0.235) (0.527,0.276) (0.535,0.315) (0.538,0.352) (0.537,0.386) (0.532,0.418) (0.523,0.446) (0.509,0.471) (0.492,0.492) (0.471,0.509) (0.446,0.523) (0.418,0.532) (0.386,0.537) (0.352,0.538) (0.315,0.535) (0.276,0.527) (0.235,0.516) (0.192,0.500) (0.148,0.481) (0.102,0.458) (0.056,0.431) (0.010,0.402) (-0.036,0.369) (-0.082,0.334) (-0.128,0.296) (-0.172,0.257) (-0.215,0.215) (-0.257,0.172) (-0.296,0.128) (-0.334,0.082) (-0.369,0.036) (-0.402,-0.010) (-0.431,-0.056) (-0.458,-0.102) (-0.481,-0.148) (-0.500,-0.192) (-0.516,-0.235) (-0.527,-0.276) (-0.535,-0.315) (-0.538,-0.352) (-0.537,-0.386) (-0.532,-0.418) (-0.523,-0.446) (-0.509,-0.471) (-0.492,-0.492) (-0.471,-0.509) (-0.446,-0.523) (-0.418,-0.532) (-0.386,-0.537) (-0.352,-0.538) (-0.315,-0.535) (-0.276,-0.527) (-0.235,-0.516) (-0.192,-0.500) (-0.148,-0.481) (-0.102,-0.458) (-0.056,-0.431) (-0.010,-0.402) (0.036,-0.369) (0.082,-0.334) (0.128,-0.296) (0.172,-0.257) (0.215,-0.215) (0.257,-0.172) (0.296,-0.128) (0.334,-0.082) (0.369,-0.036) (0.402,0.010) (0.431,0.056) (0.458,0.102) (0.481,0.148)};
\draw[couleurN] plot[smooth cycle] coordinates {(0.750,0.273) (0.773,0.337) (0.791,0.400) (0.802,0.459) (0.806,0.516) (0.804,0.568) (0.796,0.617) (0.781,0.661) (0.760,0.699) (0.733,0.733) (0.699,0.760) (0.661,0.781) (0.617,0.796) (0.568,0.804) (0.516,0.806) (0.459,0.802) (0.400,0.791) (0.337,0.773) (0.273,0.750) (0.206,0.721) (0.139,0.687) (0.071,0.647) (0.003,0.603) (-0.066,0.555) (-0.133,0.503) (-0.200,0.448) (-0.265,0.389) (-0.328,0.328) (-0.389,0.265) (-0.448,0.200) (-0.503,0.133) (-0.555,0.066) (-0.603,-0.003) (-0.647,-0.071) (-0.687,-0.139) (-0.721,-0.206) (-0.750,-0.273) (-0.773,-0.337) (-0.791,-0.400) (-0.802,-0.459) (-0.806,-0.516) (-0.804,-0.568) (-0.796,-0.617) (-0.781,-0.661) (-0.760,-0.699) (-0.733,-0.733) (-0.699,-0.760) (-0.661,-0.781) (-0.617,-0.796) (-0.568,-0.804) (-0.516,-0.806) (-0.459,-0.802) (-0.400,-0.791) (-0.337,-0.773) (-0.273,-0.750) (-0.206,-0.721) (-0.139,-0.687) (-0.071,-0.647) (-0.003,-0.603) (0.066,-0.555) (0.133,-0.503) (0.200,-0.448) (0.265,-0.389) (0.328,-0.328) (0.389,-0.265) (0.448,-0.200) (0.503,-0.133) (0.555,-0.066) (0.603,0.003) (0.647,0.071) (0.687,0.139) (0.721,0.206)};
\draw[couleurN,thick] plot[smooth cycle] coordinates {(1.000,0.337) (1.031,0.423) (1.054,0.507) (1.068,0.588) (1.074,0.665) (1.070,0.738) (1.058,0.805) (1.036,0.866) (1.006,0.920) (0.967,0.967) (0.920,1.006) (0.866,1.036) (0.805,1.058) (0.738,1.070) (0.665,1.074) (0.588,1.068) (0.507,1.054) (0.423,1.031) (0.337,1.000) (0.249,0.961) (0.160,0.916) (0.070,0.864) (-0.019,0.806) (-0.108,0.742) (-0.195,0.674) (-0.281,0.602) (-0.365,0.526) (-0.447,0.447) (-0.526,0.365) (-0.602,0.281) (-0.674,0.195) (-0.742,0.108) (-0.806,0.019) (-0.864,-0.070) (-0.916,-0.160) (-0.961,-0.249) (-1.000,-0.337) (-1.031,-0.423) (-1.054,-0.507) (-1.068,-0.588) (-1.074,-0.665) (-1.070,-0.738) (-1.058,-0.805) (-1.036,-0.866) (-1.006,-0.920) (-0.967,-0.967) (-0.920,-1.006) (-0.866,-1.036) (-0.805,-1.058) (-0.738,-1.070) (-0.665,-1.074) (-0.588,-1.068) (-0.507,-1.054) (-0.423,-1.031) (-0.337,-1.000) (-0.249,-0.961) (-0.160,-0.916) (-0.070,-0.864) (0.019,-0.806) (0.108,-0.742) (0.195,-0.674) (0.281,-0.602) (0.365,-0.526) (0.447,-0.447) (0.526,-0.365) (0.602,-0.281) (0.674,-0.195) (0.742,-0.108) (0.806,-0.019) (0.864,0.070) (0.916,0.160) (0.961,0.249)};
\draw[couleurN] plot[smooth] coordinates {(0.000,0.000) (0.100,0.040) (0.200,0.079) (0.300,0.118) (0.400,0.156) (0.500,0.192) (0.600,0.226) (0.700,0.258) (0.800,0.287) (0.900,0.313) (1.000,0.337)};
\draw[couleurN] plot[smooth] coordinates {(0.000,0.000) (0.107,0.085) (0.213,0.169) (0.320,0.253) (0.426,0.336) (0.532,0.418) (0.638,0.499) (0.743,0.578) (0.849,0.655) (0.953,0.731) (1.058,0.805)};
\draw[couleurN] plot[smooth] coordinates {(0.000,0.000) (0.085,0.107) (0.169,0.213) (0.253,0.320) (0.336,0.426) (0.418,0.532) (0.499,0.638) (0.578,0.743) (0.655,0.849) (0.731,0.953) (0.805,1.058)};
\draw[couleurN] plot[smooth] coordinates {(0.000,0.000) (0.040,0.100) (0.079,0.200) (0.118,0.300) (0.156,0.400) (0.192,0.500) (0.226,0.600) (0.258,0.700) (0.287,0.800) (0.313,0.900) (0.337,1.000)};
\draw[couleurN] plot[smooth] coordinates {(0.000,0.000) (-0.015,0.067) (-0.031,0.133) (-0.047,0.200) (-0.064,0.267) (-0.082,0.334) (-0.101,0.401) (-0.122,0.469) (-0.145,0.537) (-0.169,0.605) (-0.195,0.674)};
\draw[couleurN] plot[smooth] coordinates {(0.000,0.000) (-0.067,0.015) (-0.133,0.031) (-0.200,0.047) (-0.267,0.064) (-0.334,0.082) (-0.401,0.101) (-0.469,0.122) (-0.537,0.145) (-0.605,0.169) (-0.674,0.195)};
\draw[couleurN] plot[smooth] coordinates {(0.000,0.000) (-0.100,-0.040) (-0.200,-0.079) (-0.300,-0.118) (-0.400,-0.156) (-0.500,-0.192) (-0.600,-0.226) (-0.700,-0.258) (-0.800,-0.287) (-0.900,-0.313) (-1.000,-0.337)};
\draw[couleurN] plot[smooth] coordinates {(-0.000,-0.000) (-0.107,-0.085) (-0.213,-0.169) (-0.320,-0.253) (-0.426,-0.336) (-0.532,-0.418) (-0.638,-0.499) (-0.743,-0.578) (-0.849,-0.655) (-0.953,-0.731) (-1.058,-0.805)};
\draw[couleurN] plot[smooth] coordinates {(-0.000,-0.000) (-0.085,-0.107) (-0.169,-0.213) (-0.253,-0.320) (-0.336,-0.426) (-0.418,-0.532) (-0.499,-0.638) (-0.578,-0.743) (-0.655,-0.849) (-0.731,-0.953) (-0.805,-1.058)};
\draw[couleurN] plot[smooth] coordinates {(-0.000,-0.000) (-0.040,-0.100) (-0.079,-0.200) (-0.118,-0.300) (-0.156,-0.400) (-0.192,-0.500) (-0.226,-0.600) (-0.258,-0.700) (-0.287,-0.800) (-0.313,-0.900) (-0.337,-1.000)};
\draw[couleurN] plot[smooth] coordinates {(0.000,0.000) (0.015,-0.067) (0.031,-0.133) (0.047,-0.200) (0.064,-0.267) (0.082,-0.334) (0.101,-0.401) (0.122,-0.469) (0.145,-0.537) (0.169,-0.605) (0.195,-0.674)};
\draw[couleurN] plot[smooth] coordinates {(0.000,0.000) (0.067,-0.015) (0.133,-0.031) (0.200,-0.047) (0.267,-0.064) (0.334,-0.082) (0.401,-0.101) (0.469,-0.122) (0.537,-0.145) (0.605,-0.169) (0.674,-0.195)};
    \draw[rougecours,densely dashed,thick] (0,0) circle (0.6);
    \fill[rougecours] (0,0) circle (0.9pt);
    \draw[rougecours,->,>=Stealth] (1.25,-1.15) node[right,font=\small,align=left] {disque de rayon\\ $1 - k = 0{,}6$} -- (0.40,-0.47);
    \node[couleurN,font=\small,anchor=south west] at (0.85,0.95) {$\psi\big(\Bf(0, 1)\big)$};
  \end{scope}
\end{tikzpicture}
\end{center}

\noindent Pour $\psi(x, y) = (x + 0{,}4\sin y, y + 0{,}4\sin x)$, la perturbation $\varphi$ est $0{,}4$-lipschitzienne pour la norme euclidienne, puisque la norme subordonnée de sa jacobienne $\begin{pmatrix} 0 & 0{,}4\cos y \\ 0{,}4\cos x & 0 \end{pmatrix}$ est au plus $0{,}4$. Le disque unité et son quadrillage polaire, en violet, sont déformés par $\psi$ en la région bleue, qui contient bien le disque de rayon $0{,}6$, en pointillés rouges : tout point de ce disque a un unique antécédent dans le disque unité.

\begin{cadre}
\propo{5} \emph{(Newton à jacobienne figée.)} Soit $E$ et $F$ deux espaces de Banach, $f : U \to F$ différentiable, et $a \in U$ tel que $\dd f(a) \in \Isom(E, F)$. On note $L = \dd f(a)^{-1}$, et l'on suppose que $r > 0$ vérifie $\Bf(a, r) \subset U$ et
\[ \nop{L \circ \big(\dd f(x) - \dd f(a)\big)} \leqslant \frac{1}{2} \qquad \text{pour tout } x \in \Bf(a, r). \]
\begin{enumerate}[label=(\roman*),topsep=3pt,itemsep=2pt]
  \item Pour tous $x$, $x' \in \Bf(a, r)$, $\rmath{\norme{x - x'} \leqslant 2\,\norme{L\cdot\big(f(x) - f(x')\big)}}$.
  \item Pour tout $y \in F$ tel que $\norme{L\cdot(y - f(a))} \leqslant \dfrac{r}{2}$, l'équation $f(x) = y$ a une unique solution $x \in \Bf(a, r)$, qui vérifie $\norme{x - a} \leqslant 2\norme{L\cdot(y - f(a))}$. C'est la limite de la suite définie par $x_{0} = a$ et
  \[ \rmath{x_{n + 1} = x_{n} - L\cdot\big(f(x_{n}) - y\big)}, \qquad \text{avec } \norme{x_{n} - x} \leqslant \frac{r}{2^{n}}. \]
\end{enumerate}
\end{cadre}

\begin{demo}
On se ramène à la proposition 4 en posant, pour $h \in \Bf(0, r)$,
\[ \psi(h) = L\cdot\big(f(a + h) - f(a)\big), \qquad \varphi(h) = \psi(h) - h. \]
Alors $\varphi(0) = 0$, et $\varphi$ est différentiable, de différentielle $\dd\varphi(h) = L \circ \dd f(a + h) - \id = L \circ (\dd f(a + h) - \dd f(a))$, de norme au plus $\dfrac{1}{2}$. La boule $\Bf(0, r)$ est convexe : par l'inégalité des accroissements finis (chapitre 2, théorème 4), $\varphi$ est $\dfrac{1}{2}$-lipschitzienne. La proposition 4 s'applique avec $k = \dfrac{1}{2}$.

(i) découle de la proposition 4 (i), puisque $\psi(h) - \psi(h') = L\cdot(f(a + h) - f(a + h'))$.

(ii) Pour $x = a + h$, on a $f(x) = y$ si et seulement si $\psi(h) = z$, avec $z = L\cdot(y - f(a))$, car $L$ est injective. Comme $\norme{z} \leqslant \dfrac{r}{2} = (1 - k)r$, la proposition 4 (ii) donne l'existence, l'unicité et la majoration $\norme{h} \leqslant 2\norme{z}$. Enfin, avec $x_{n} = a + h_{n}$, la relation $h_{n + 1} = z - \varphi(h_{n})$ s'écrit
\[ x_{n + 1} - a = L\cdot(y - f(a)) - L\cdot(f(x_{n}) - f(a)) + (x_{n} - a), \]
c'est-à-dire $x_{n + 1} = x_{n} - L\cdot(f(x_{n}) - y)$ : c'est la méthode à jacobienne figée, et $\norme{x_{n} - x} = \norme{h_{n} - h} \leqslant 2^{-n}r$.
\end{demo}

\begin{remarques}
  \item La convergence est géométrique, et non quadratique : chaque étape gagne un nombre de décimales à peu près constant, ce qu'on observe sur la figure précédente. En contrepartie, l'hypothèse ne demande que la continuité de $\dd f$ en $a$, qui assure l'existence d'un tel $r$ : c'est exactement l'hypothèse $\Cc^{1}$ du théorème d'inversion locale.
  \item La proposition est \emph{quantitative} : elle donne une boule où la solution existe, une boule où elle est unique, et une majoration de l'erreur commise après $n$ étapes. Le théorème d'inversion locale qui suit n'en retient que l'existence des voisinages, mais un calcul effectif repose sur ces estimations.
\end{remarques}
%===============================================================================
\section[Inversion locale]{Le théorème d'inversion locale}
%===============================================================================

\subsection{Difféomorphismes}

Une application linéaire de $\R^{p}$ dans lui-même est bijective dès que sa matrice est inversible, et sa réciproque est alors linéaire. Le théorème d'inversion locale transporte cette propriété aux applications différentiables, mais seulement \emph{près d'un point} : là où la différentielle est inversible, l'application est inversible, et sa réciproque est aussi régulière qu'elle.

\begin{center}
\begin{tabularx}{15cm}{||X}
\noindent \textbf{\underline{Définition :}} Soit $V$ un ouvert de $E$ et $W$ un ouvert de $F$. Une application $f : V \to W$ est un \mot{$\Cc^{1}$-difféomorphisme} de $V$ sur $W$ si elle est bijective, de classe $\Cc^{1}$ sur $V$, et si sa réciproque $f^{-1} : W \to V$ est de classe $\Cc^{1}$ sur $W$.

\noindent Une application $f : U \to F$ est un \mot{$\Cc^{1}$-difféomorphisme local au voisinage de $a \in U$} s'il existe un ouvert $V \subset U$ contenant $a$ et un ouvert $W$ de $F$ contenant $f(a)$ tels que la restriction de $f$ à $V$ soit un $\Cc^{1}$-difféomorphisme de $V$ sur $W$.
\end{tabularx}
\end{center}

\begin{cadre}
\propo{6} Si $f : V \to W$ est un $\Cc^{1}$-difféomorphisme, alors, pour tout $x \in V$, $\dd f(x) \in \Isom(E, F)$ et
\[ \rmath{\dd(f^{-1})\big(f(x)\big) = \dd f(x)^{-1}}. \]
En particulier, s'il existe un $\Cc^{1}$-difféomorphisme entre un ouvert non vide de $\R^{p}$ et un ouvert de $\R^{n}$, alors $p = n$.
\end{cadre}

\begin{demo}
Soit $g = f^{-1}$ et $y = f(x)$. La règle de la chaîne (chapitre 2, théorème 3), appliquée à $g \circ f = \id_{V}$ et à $f \circ g = \id_{W}$, donne $\dd g(y) \circ \dd f(x) = \id_{E}$ et $\dd f(x) \circ \dd g(y) = \id_{F}$. Ainsi $\dd f(x)$ est bijective, de réciproque $\dd g(y)$ continue. Entre $\R^{p}$ et $\R^{n}$, une application linéaire bijective impose $p = n$.
\end{demo}

\begin{remarques}
  \item \emph{Piège.} Une bijection de classe $\Cc^{1}$ n'est pas forcément un difféomorphisme. La fonction $x \mapsto x^{3}$ est une bijection de classe $\Cc^{1}$ de $\R$ sur $\R$, mais sa réciproque $x \mapsto x^{1/3}$, prolongée par continuité, n'est pas dérivable en $0$ : c'est la proposition 6, puisque la dérivée de $x^{3}$ s'annule en $0$. L'hypothèse «~différentielle inversible~» du théorème qui suit ne peut donc pas être omise.
  \item Pour une fonction $f : I \to \R$ de classe $\Cc^{1}$ sur un intervalle ouvert, dont la dérivée ne s'annule pas, on sait que $f$ est strictement monotone, que $f(I)$ est un intervalle ouvert et que $f^{-1}$ est de classe $\Cc^{1}$ (\emph{Rappel.}). En dimension supérieure, aucun argument de monotonie n'est disponible : c'est le point fixe qui le remplace.
\end{remarques}

\subsection{L'énoncé et sa démonstration}

\begin{center}
\begin{tabularx}{15cm}{||X||}
\hline \\
\textbf{\underline{\textcolor{rougecours}{Théorème 1 : Inversion locale :}}} \\
Soit $E$ et $F$ deux espaces de Banach, $U$ un ouvert de $E$, $f : U \to F$ de classe $\Cc^{1}$, et $a \in U$ tel que $\dd f(a) \in \Isom(E, F)$. Il existe un ouvert $V \subset U$ contenant $a$ et un ouvert $W$ de $F$ contenant $f(a)$ tels que $f$ soit un $\Cc^{1}$-difféomorphisme de $V$ sur $W$. De plus, pour tout $y \in W$,
\[ \rmath{\dd(f^{-1})(y) = \dd f\big(f^{-1}(y)\big)^{-1}}. \]
\medbreak
\\ \hline
\end{tabularx}
\end{center}

\begin{demo}
On note $L = \dd f(a)^{-1}$.

\etape{Choix du rayon.} L'application $x \mapsto L \circ (\dd f(x) - \dd f(a))$ est continue, puisque $\dd f$ l'est, et nulle en $a$. Il existe donc $r > 0$ tel que $\Bf(a, r) \subset U$ et $\nop{L \circ (\dd f(x) - \dd f(a))} \leqslant \dfrac{1}{2}$ pour tout $x \in \Bf(a, r)$ : les hypothèses de la proposition 5 sont satisfaites.

\etape{Les ouverts $V$ et $W$.} On pose
\[ W = \left\{y \in F \;\Big|\; \norme{L\cdot(y - f(a))} < \frac{r}{2}\right\}, \qquad V = \Bo(a, r) \cap f^{-1}(W). \]
L'ensemble $W$ est l'image réciproque d'un ouvert de $\R$ par l'application continue $y \mapsto \norme{L\cdot(y - f(a))}$ : c'est un ouvert contenant $f(a)$. L'ensemble $V$ est ouvert, puisque $f$ est continue, et contient $a$. Pour $y \in W$, la proposition 5 (ii) fournit une solution $x \in \Bf(a, r)$ de $f(x) = y$, avec $\norme{x - a} \leqslant 2\norme{L\cdot(y - f(a))} < r$ : ainsi $x \in V$, et $f : V \to W$ est surjective. Elle est injective par la proposition 5 (i). Notons $g : W \to V$ sa réciproque.

\etape{Continuité de $g$.} Pour $y$, $y' \in W$, la proposition 5 (i), appliquée à $x = g(y)$ et $x' = g(y')$, donne
\[ \norme{g(y) - g(y')} \leqslant 2\nop{L}\,\norme{y - y'} : \]
$g$ est lipschitzienne.

\etape{Inversibilité de $\dd f$ sur $V$.} Pour $x \in V$, on écrit $\dd f(x) = \dd f(a) \circ \big(\id_{E} + L \circ (\dd f(x) - \dd f(a))\big)$. Le second facteur est de la forme $\id_{E} + u$ avec $\nop{u} \leqslant \dfrac{1}{2} < 1$ : il appartient à $\GL(E)$ (chapitre 1, théorème 12). Ainsi $\dd f(x) \in \Isom(E, F)$.

\etape{Différentiabilité de $g$.} Soit $y_{0} \in W$, $x_{0} = g(y_{0})$, et, pour $y_{0} + k \in W$, $h = g(y_{0} + k) - x_{0}$. La différentiabilité de $f$ en $x_{0}$ s'écrit $f(x_{0} + h) - f(x_{0}) = \dd f(x_{0})\cdot h + \norme{h}\,\varepsilon(h)$, avec $\varepsilon(h) \to 0$ quand $h \to 0$, soit
\[ k = \dd f(x_{0})\cdot h + \norme{h}\,\varepsilon(h), \qquad h = \dd f(x_{0})^{-1}\cdot k - \norme{h}\,\dd f(x_{0})^{-1}\cdot\varepsilon(h). \]
Comme $\norme{h} \leqslant 2\nop{L}\,\norme{k}$ par l'étape 3, et que $h \to 0$ quand $k \to 0$, le dernier terme est un $o(k)$ : $g$ est différentiable en $y_{0}$, de différentielle $\dd f(x_{0})^{-1}$.

\etape{Classe $\Cc^{1}$.} On a $\dd g = \mathrm{inv} \circ \dd f \circ g$, où $\mathrm{inv} : u \mapsto u^{-1}$ est continue sur $\Isom(E, F)$ (proposition 2). Par composition d'applications continues, $\dd g$ est continue.
\end{demo}

\begin{remarques}
  \item La démonstration est constructive : pour $y \in W$, l'antécédent $f^{-1}(y)$ est la limite de la suite de Newton à jacobienne figée $x_{n + 1} = x_{n} - \dd f(a)^{-1}\cdot(f(x_{n}) - y)$, issue de $a$. Calculer une réciproque locale, c'est résoudre une équation, et la méthode de Newton le fait.
  \item Le théorème est \emph{local} : même si $\dd f(x)$ est inversible en tout point, $f$ peut ne pas être injective sur $U$, comme le montre l'exemple de $z \mapsto z^{2}$ ci-dessous.
  \item Si $f$ est de classe $\Cc^{k}$, la réciproque l'est aussi, puisque $\dd g = \mathrm{inv} \circ \dd f \circ g$ et que l'inversion est elle-même de classe $\Cc^{\infty}$ : la démonstration, par récurrence sur $k$, demande les différentielles d'ordre supérieur du chapitre 4.
\end{remarques}

\begin{cadre}
\propo{7} \emph{(Inversion globale.)} Soit $E$ et $F$ deux espaces de Banach et $f : U \to F$ de classe $\Cc^{1}$, \emph{injective}, telle que $\dd f(x) \in \Isom(E, F)$ pour tout $x \in U$. Alors $f(U)$ est un ouvert de $F$, et $f$ est un $\Cc^{1}$-difféomorphisme de $U$ sur $f(U)$.
\end{cadre}

\begin{demo}
Soit $y = f(x) \in f(U)$. Le théorème 1, appliqué en $x$, fournit des ouverts $V_{x} \ni x$ et $W_{x} \ni y$ avec $W_{x} = f(V_{x}) \subset f(U)$ : $f(U)$ est voisinage de chacun de ses points, donc ouvert. La réciproque $g : f(U) \to U$ existe, puisque $f$ est injective, et elle coïncide sur $W_{x}$ avec la réciproque locale fournie par le théorème 1, de classe $\Cc^{1}$. Être de classe $\Cc^{1}$ est une propriété locale : $g$ est de classe $\Cc^{1}$ sur $f(U)$.
\end{demo}

\newpage
\subsection{Exemples}

\begin{exemples}
  \item \emph{Coordonnées polaires.} L'application $P(r, \theta) = (r\cos\theta, r\sin\theta)$ a pour jacobien $r$ (chapitre 2, section 2.4) : elle est un difféomorphisme local au voisinage de tout point où $r \neq 0$. Elle n'est pas injective sur $\left]0, +\infty\right[ \times \R$, puisqu'elle est $2\pi$-périodique en $\theta$, mais elle l'est sur $U = \left]0, +\infty\right[ \times \left]-\pi, \pi\right[$. Par la proposition 7, elle est un $\Cc^{1}$-difféomorphisme de $U$ sur le plan privé de la demi-droite $\R_{-} \times \{0\}$, et sa réciproque fournit l'angle polaire, dont le chapitre 2 a calculé le gradient.
  \item \emph{L'élévation au carré dans $\C$.} Identifions $\C$ à $\R^{2}$ et posons $f(x, y) = (x^{2} - y^{2}, 2xy)$, qui correspond à $z \mapsto z^{2}$. Sa jacobienne est $\begin{pmatrix} 2x & -2y \\ 2y & 2x \end{pmatrix}$, de déterminant $4(x^{2} + y^{2})$ : $f$ est un difféomorphisme local au voisinage de tout point non nul. Pourtant $f(-z) = f(z)$ : $f$ n'est injective sur aucun ouvert contenant à la fois $z$ et $-z$.
\end{exemples}

% ---- f(z) = z^2 ; a = (1, 0.5), carres de demi-cote 0.25 centres en a et -a ; f(a) = (0.75, 1) ; jacobien en a : 5 (calcules)
\begin{center}
\begin{tikzpicture}[scale=1.9]
  \begin{scope}
    \draw[gray,->] (-1.6,0) -- (1.6,0) node[below,font=\small] {$x$};
    \draw[gray,->] (0,-1.0) -- (0,1.0) node[left,font=\small] {$y$};
\draw[couleurQ] plot coordinates {(0.750,0.250) (0.750,0.312) (0.750,0.375) (0.750,0.438) (0.750,0.500) (0.750,0.562) (0.750,0.625) (0.750,0.688) (0.750,0.750)};
\draw[couleurQ] plot coordinates {(0.750,0.250) (0.812,0.250) (0.875,0.250) (0.938,0.250) (1.000,0.250) (1.062,0.250) (1.125,0.250) (1.188,0.250) (1.250,0.250)};
\draw[couleurQ] plot coordinates {(0.917,0.250) (0.917,0.312) (0.917,0.375) (0.917,0.438) (0.917,0.500) (0.917,0.562) (0.917,0.625) (0.917,0.688) (0.917,0.750)};
\draw[couleurQ] plot coordinates {(0.750,0.417) (0.812,0.417) (0.875,0.417) (0.938,0.417) (1.000,0.417) (1.062,0.417) (1.125,0.417) (1.188,0.417) (1.250,0.417)};
\draw[couleurQ] plot coordinates {(1.083,0.250) (1.083,0.312) (1.083,0.375) (1.083,0.438) (1.083,0.500) (1.083,0.562) (1.083,0.625) (1.083,0.688) (1.083,0.750)};
\draw[couleurQ] plot coordinates {(0.750,0.583) (0.812,0.583) (0.875,0.583) (0.938,0.583) (1.000,0.583) (1.062,0.583) (1.125,0.583) (1.188,0.583) (1.250,0.583)};
\draw[couleurQ] plot coordinates {(1.250,0.250) (1.250,0.312) (1.250,0.375) (1.250,0.438) (1.250,0.500) (1.250,0.562) (1.250,0.625) (1.250,0.688) (1.250,0.750)};
\draw[couleurQ] plot coordinates {(0.750,0.750) (0.812,0.750) (0.875,0.750) (0.938,0.750) (1.000,0.750) (1.062,0.750) (1.125,0.750) (1.188,0.750) (1.250,0.750)};
\draw[couleurQ] plot coordinates {(-1.250,-0.750) (-1.250,-0.688) (-1.250,-0.625) (-1.250,-0.562) (-1.250,-0.500) (-1.250,-0.438) (-1.250,-0.375) (-1.250,-0.312) (-1.250,-0.250)};
\draw[couleurQ] plot coordinates {(-1.250,-0.750) (-1.188,-0.750) (-1.125,-0.750) (-1.062,-0.750) (-1.000,-0.750) (-0.938,-0.750) (-0.875,-0.750) (-0.812,-0.750) (-0.750,-0.750)};
\draw[couleurQ] plot coordinates {(-1.083,-0.750) (-1.083,-0.688) (-1.083,-0.625) (-1.083,-0.562) (-1.083,-0.500) (-1.083,-0.438) (-1.083,-0.375) (-1.083,-0.312) (-1.083,-0.250)};
\draw[couleurQ] plot coordinates {(-1.250,-0.583) (-1.188,-0.583) (-1.125,-0.583) (-1.062,-0.583) (-1.000,-0.583) (-0.938,-0.583) (-0.875,-0.583) (-0.812,-0.583) (-0.750,-0.583)};
\draw[couleurQ] plot coordinates {(-0.917,-0.750) (-0.917,-0.688) (-0.917,-0.625) (-0.917,-0.562) (-0.917,-0.500) (-0.917,-0.438) (-0.917,-0.375) (-0.917,-0.312) (-0.917,-0.250)};
\draw[couleurQ] plot coordinates {(-1.250,-0.417) (-1.188,-0.417) (-1.125,-0.417) (-1.062,-0.417) (-1.000,-0.417) (-0.938,-0.417) (-0.875,-0.417) (-0.812,-0.417) (-0.750,-0.417)};
\draw[couleurQ] plot coordinates {(-0.750,-0.750) (-0.750,-0.688) (-0.750,-0.625) (-0.750,-0.562) (-0.750,-0.500) (-0.750,-0.438) (-0.750,-0.375) (-0.750,-0.312) (-0.750,-0.250)};
\draw[couleurQ] plot coordinates {(-1.250,-0.250) (-1.188,-0.250) (-1.125,-0.250) (-1.062,-0.250) (-1.000,-0.250) (-0.938,-0.250) (-0.875,-0.250) (-0.812,-0.250) (-0.750,-0.250)};
    \fill[rougecours] (1,0.5) circle (1.1pt);
    \node[rougecours,font=\small,anchor=south west] at (1.26,0.76) {$a$};
    \fill[rougecours] (-1,-0.5) circle (1.1pt);
    \node[rougecours,font=\small,anchor=north east] at (-1.26,-0.76) {$-a$};
    \fill[gray] (0,0) circle (0.8pt);
  \end{scope}
  \draw[gray,->,>=Stealth] (1.7,0.2) to[bend left=20] node[above,font=\small] {$z \mapsto z^{2}$} (2.7,0.2);
  \begin{scope}[xshift=3.35cm]
    \draw[gray,->] (-0.5,0) -- (1.6,0) node[below,font=\small] {$X$};
    \draw[gray,->] (0,-1.0) -- (0,2.05) node[left,font=\small] {$Y$};
\draw[couleurN] plot[smooth] coordinates {(0.500,0.375) (0.465,0.469) (0.422,0.562) (0.371,0.656) (0.312,0.750) (0.246,0.844) (0.172,0.938) (0.090,1.031) (0.000,1.125)};
\draw[couleurN] plot[smooth] coordinates {(0.500,0.375) (0.598,0.406) (0.703,0.438) (0.816,0.469) (0.938,0.500) (1.066,0.531) (1.203,0.562) (1.348,0.594) (1.500,0.625)};
\draw[couleurN] plot[smooth] coordinates {(0.778,0.458) (0.743,0.573) (0.700,0.688) (0.649,0.802) (0.590,0.917) (0.524,1.031) (0.450,1.146) (0.368,1.260) (0.278,1.375)};
\draw[couleurN] plot[smooth] coordinates {(0.389,0.625) (0.487,0.677) (0.592,0.729) (0.705,0.781) (0.826,0.833) (0.955,0.885) (1.092,0.937) (1.237,0.990) (1.389,1.042)};
\draw[couleurN] plot[smooth] coordinates {(1.111,0.542) (1.076,0.677) (1.033,0.812) (0.982,0.948) (0.924,1.083) (0.857,1.219) (0.783,1.354) (0.701,1.490) (0.611,1.625)};
\draw[couleurN] plot[smooth] coordinates {(0.222,0.875) (0.320,0.948) (0.425,1.021) (0.539,1.094) (0.660,1.167) (0.789,1.240) (0.925,1.312) (1.070,1.385) (1.222,1.458)};
\draw[couleurN] plot[smooth] coordinates {(1.500,0.625) (1.465,0.781) (1.422,0.938) (1.371,1.094) (1.312,1.250) (1.246,1.406) (1.172,1.562) (1.090,1.719) (1.000,1.875)};
\draw[couleurN] plot[smooth] coordinates {(0.000,1.125) (0.098,1.219) (0.203,1.312) (0.316,1.406) (0.438,1.500) (0.566,1.594) (0.703,1.688) (0.848,1.781) (1.000,1.875)};
    \fill[rougecours] (0.75,1) circle (1.1pt);
    \draw[rougecours,->,>=Stealth] (1.75,0.3) node[below,font=\small] {$f(a) = f(-a)$} -- (0.79,0.97);
    \fill[gray] (0,0) circle (0.8pt);
  \end{scope}
\end{tikzpicture}
\end{center}

\noindent Les deux petits carrés violets, centrés en $a = (1 ; 0{,}5)$ et en $-a$, ont la même image bleue, un quadrilatère curviligne autour de $f(a) = (0{,}75 ; 1)$ dont l'aire vaut environ $5$ fois celle d'un carré, le jacobien en $a$. Restreinte à chacun des deux carrés, $f$ est bijective, et c'est tout ce qu'affirme le théorème : sur un domaine qui contient les deux carrés, elle ne l'est plus.

\begin{exemple}
\emph{Racine carrée d'une matrice.} Soit $f : \Mn{n} \to \Mn{n}$, $f(X) = X^{2}$, avec $\K = \R$. Elle est de classe $\Cc^{1}$, de différentielle $\dd f(X)\cdot H = XH + HX$ (chapitre 2), d'où $\dd f(I_{n})\cdot H = 2H$ : $\dd f(I_{n}) = 2\,\id$ est inversible. Par le théorème 1, il existe des voisinages $V$ de $I_{n}$ et $W$ de $I_{n}$ tels que toute matrice $B \in W$ ait une unique «~racine carrée~» $X \in V$, qui dépend de $B$ de façon $\Cc^{1}$. La méthode à jacobienne figée s'écrit
\[ X_{k + 1} = X_{k} - \frac{1}{2}\big(X_{k}^{2} - B\big), \qquad X_{0} = I_{n}. \]
Pour la matrice $B$ ci-dessous, on obtient la première itérée $X_{1}$, et la suite converge vers la limite $X$ :
\[ B = \begin{pmatrix} 1{,}2 & 0{,}3 \\ 0{,}1 & 0{,}9 \end{pmatrix}, \qquad X_{1} = \begin{pmatrix} 1{,}1 & 0{,}15 \\ 0{,}05 & 0{,}95 \end{pmatrix}, \qquad X \approx \begin{pmatrix} 1{,}092\,14 & 0{,}147\,28 \\ 0{,}049\,09 & 0{,}944\,87 \end{pmatrix}. \]
Les écarts $X_{k} - X$, mesurés par le plus grand coefficient en valeur absolue, valent environ $7{,}9 \times 10^{-3}$, $1{,}0 \times 10^{-3}$, $1{,}3 \times 10^{-4}$, $1{,}6 \times 10^{-5}$ : chaque étape gagne à peu près un facteur $8$. On vérifie numériquement que la limite a bien pour carré $B$.
\end{exemple}

\begin{remarque}
\emph{Piège.} L'hypothèse «~de classe $\Cc^{1}$~» ne peut pas être affaiblie en «~différentiable~». La fonction $f(x) = \dfrac{x}{2} + x^{2}\sin\dfrac{1}{x}$, prolongée par $f(0) = 0$, est dérivable sur $\R$, avec $f'(0) = \dfrac{1}{2} \neq 0$. Pourtant $f'(x) = \dfrac{1}{2} + 2x\sin\dfrac{1}{x} - \cos\dfrac{1}{x}$ vaut $-\dfrac{1}{2}$ aux points $\dfrac{1}{2k\pi}$ et $\dfrac{3}{2}$ aux points $\dfrac{1}{(2k + 1)\pi}$ : la fonction n'est monotone sur aucun intervalle contenant $0$, donc injective sur aucun voisinage de $0$, puisqu'une fonction continue et injective sur un intervalle y est strictement monotone. Sa dérivée n'est pas continue en $0$, comme celle de l'exemple $x^{2}\sin\dfrac{1}{x}$ du chapitre 2.
\end{remarque}

% ---- f(x) = x/2 + x^2 sin(1/x) sur ]0 ; 0.1] ; f'(1/(2k pi)) = -1/2 (calcule)
\begin{center}
\begin{tikzpicture}
  \begin{axis}[width=13cm,height=6.2cm,xmin=0,xmax=0.102,ymin=-0.003,ymax=0.062,axis lines=middle,
    tick label style={font=\scriptsize,/pgf/number format/fixed,/pgf/number format/use comma},
    xtick={0.02,0.04,0.06,0.08,0.1},ytick={0.02,0.04,0.06},scaled ticks=false,clip=false]
    \addplot[couleurN,thick] coordinates {(0.00040,0.000200) (0.00042,0.000211) (0.00044,0.000222) (0.00047,0.000233) (0.00049,0.000243) (0.00051,0.000254) (0.00053,0.000265) (0.00055,0.000276) (0.00057,0.000287) (0.00060,0.000298) (0.00062,0.000309) (0.00064,0.000320) (0.00066,0.000331) (0.00068,0.000341) (0.00071,0.000352) (0.00073,0.000363) (0.00075,0.000374) (0.00077,0.000385) (0.00079,0.000396) (0.00081,0.000407) (0.00084,0.000419) (0.00086,0.000429) (0.00088,0.000439) (0.00090,0.000450) (0.00092,0.000462) (0.00095,0.000473) (0.00097,0.000483) (0.00099,0.000494) (0.00101,0.000505) (0.00103,0.000517) (0.00105,0.000527) (0.00108,0.000537) (0.00110,0.000549) (0.00112,0.000561) (0.00114,0.000571) (0.00116,0.000580) (0.00118,0.000594) (0.00121,0.000602) (0.00123,0.000614) (0.00125,0.000627) (0.00127,0.000637) (0.00129,0.000647) (0.00132,0.000657) (0.00134,0.000669) (0.00136,0.000681) (0.00138,0.000692) (0.00140,0.000702) (0.00142,0.000710) (0.00145,0.000724) (0.00147,0.000735) (0.00149,0.000743) (0.00151,0.000758) (0.00153,0.000765) (0.00156,0.000780) (0.00158,0.000787) (0.00160,0.000799) (0.00162,0.000813) (0.00164,0.000820) (0.00166,0.000830) (0.00169,0.000845) (0.00171,0.000857) (0.00173,0.000865) (0.00175,0.000873) (0.00177,0.000884) (0.00180,0.000895) (0.00182,0.000907) (0.00184,0.000918) (0.00186,0.000930) (0.00188,0.000940) (0.00190,0.000951) (0.00193,0.000960) (0.00195,0.000970) (0.00197,0.000981) (0.00199,0.000994) (0.00201,0.001008) (0.00204,0.001022) (0.00206,0.001031) (0.00208,0.001038) (0.00210,0.001046) (0.00212,0.001061) (0.00214,0.001077) (0.00217,0.001084) (0.00219,0.001089) (0.00221,0.001106) (0.00223,0.001120) (0.00225,0.001123) (0.00227,0.001136) (0.00230,0.001153) (0.00232,0.001155) (0.00234,0.001170) (0.00236,0.001185) (0.00238,0.001186) (0.00241,0.001208) (0.00243,0.001212) (0.00245,0.001224) (0.00247,0.001239) (0.00249,0.001241) (0.00251,0.001264) (0.00254,0.001262) (0.00256,0.001286) (0.00258,0.001284) (0.00260,0.001307) (0.00262,0.001306) (0.00265,0.001329) (0.00267,0.001328) (0.00269,0.001351) (0.00271,0.001348) (0.00273,0.001374) (0.00275,0.001370) (0.00278,0.001395) (0.00280,0.001394) (0.00282,0.001413) (0.00284,0.001421) (0.00286,0.001428) (0.00289,0.001450) (0.00291,0.001445) (0.00293,0.001472) (0.00295,0.001472) (0.00297,0.001484) (0.00299,0.001504) (0.00302,0.001499) (0.00304,0.001525) (0.00306,0.001531) (0.00308,0.001533) (0.00310,0.001561) (0.00313,0.001558) (0.00315,0.001569) (0.00317,0.001594) (0.00319,0.001588) (0.00321,0.001603) (0.00323,0.001627) (0.00326,0.001621) (0.00328,0.001635) (0.00330,0.001661) (0.00332,0.001655) (0.00334,0.001665) (0.00337,0.001693) (0.00339,0.001693) (0.00341,0.001694) (0.00343,0.001722) (0.00345,0.001733) (0.00347,0.001726) (0.00350,0.001746) (0.00352,0.001771) (0.00354,0.001767) (0.00356,0.001769) (0.00358,0.001798) (0.00360,0.001813) (0.00363,0.001805) (0.00365,0.001815) (0.00367,0.001845) (0.00369,0.001855) (0.00371,0.001846) (0.00374,0.001859) (0.00376,0.001890) (0.00378,0.001899) (0.00380,0.001890) (0.00382,0.001901) (0.00384,0.001931) (0.00387,0.001946) (0.00389,0.001938) (0.00391,0.001940) (0.00393,0.001968) (0.00395,0.001992) (0.00398,0.001991) (0.00400,0.001984) (0.00402,0.002000) (0.00404,0.002031) (0.00406,0.002046) (0.00408,0.002039) (0.00411,0.002036) (0.00413,0.002058) (0.00415,0.002089) (0.00417,0.002100) (0.00419,0.002092) (0.00422,0.002090) (0.00424,0.002112) (0.00426,0.002143) (0.00428,0.002157) (0.00430,0.002150) (0.00432,0.002145) (0.00435,0.002160) (0.00437,0.002191) (0.00439,0.002214) (0.00441,0.002215) (0.00443,0.002205) (0.00446,0.002208) (0.00448,0.002232) (0.00450,0.002263) (0.00452,0.002280) (0.00454,0.002276) (0.00456,0.002267) (0.00459,0.002273) (0.00461,0.002299) (0.00463,0.002330) (0.00465,0.002347) (0.00467,0.002344) (0.00469,0.002334) (0.00472,0.002336) (0.00474,0.002358) (0.00476,0.002389) (0.00478,0.002414) (0.00480,0.002419) (0.00483,0.002410) (0.00485,0.002403) (0.00487,0.002413) (0.00489,0.002440) (0.00491,0.002471) (0.00493,0.002492) (0.00496,0.002494) (0.00498,0.002484) (0.00500,0.002478) (0.00502,0.002487) (0.00504,0.002513) (0.00507,0.002545) (0.00509,0.002569) (0.00511,0.002576) (0.00513,0.002569) (0.00515,0.002559) (0.00517,0.002561) (0.00520,0.002579) (0.00522,0.002609) (0.00524,0.002640) (0.00526,0.002659) (0.00528,0.002661) (0.00531,0.002652) (0.00533,0.002644) (0.00535,0.002646) (0.00537,0.002664) (0.00539,0.002694) (0.00541,0.002725) (0.00544,0.002747) (0.00546,0.002754) (0.00548,0.002748) (0.00550,0.002738) (0.00552,0.002734) (0.00555,0.002743) (0.00557,0.002767) (0.00559,0.002799) (0.00561,0.002829) (0.00563,0.002848) (0.00565,0.002853) (0.00568,0.002846) (0.00570,0.002836) (0.00572,0.002831) (0.00574,0.002838) (0.00576,0.002860) (0.00579,0.002890) (0.00581,0.002922) (0.00583,0.002946) (0.00585,0.002958) (0.00587,0.002957) (0.00589,0.002948) (0.00592,0.002938) (0.00594,0.002936) (0.00596,0.002946) (0.00598,0.002968) (0.00600,0.002999) (0.00602,0.003031) (0.00605,0.003056) (0.00607,0.003071) (0.00609,0.003073) (0.00611,0.003065) (0.00613,0.003055) (0.00616,0.003048) (0.00618,0.003051) (0.00620,0.003066) (0.00622,0.003091) (0.00624,0.003123) (0.00626,0.003154) (0.00629,0.003179) (0.00631,0.003194) (0.00633,0.003196) (0.00635,0.003190) (0.00637,0.003179) (0.00640,0.003171) (0.00642,0.003170) (0.00644,0.003179) (0.00646,0.003199) (0.00648,0.003228) (0.00650,0.003261) (0.00653,0.003291) (0.00655,0.003314) (0.00657,0.003328) (0.00659,0.003330) (0.00661,0.003324) (0.00664,0.003314) (0.00666,0.003304) (0.00668,0.003300) (0.00670,0.003305) (0.00672,0.003321) (0.00674,0.003346) (0.00677,0.003377) (0.00679,0.003409) (0.00681,0.003438) (0.00683,0.003460) (0.00685,0.003473) (0.00688,0.003476) (0.00690,0.003470) (0.00692,0.003460) (0.00694,0.003450) (0.00696,0.003444) (0.00698,0.003445) (0.00701,0.003455) (0.00703,0.003475) (0.00705,0.003502) (0.00707,0.003534) (0.00709,0.003566) (0.00712,0.003595) (0.00714,0.003617) (0.00716,0.003630) (0.00718,0.003635) (0.00720,0.003631) (0.00722,0.003622) (0.00725,0.003611) (0.00727,0.003603) (0.00729,0.003599) (0.00731,0.003603) (0.00733,0.003615) (0.00735,0.003636) (0.00738,0.003663) (0.00740,0.003695) (0.00742,0.003728) (0.00744,0.003757) (0.00746,0.003782) (0.00749,0.003799) (0.00751,0.003807) (0.00753,0.003808) (0.00755,0.003802) (0.00757,0.003792) (0.00759,0.003782) (0.00762,0.003773) (0.00764,0.003769) (0.00766,0.003772) (0.00768,0.003783) (0.00770,0.003802) (0.00773,0.003827) (0.00775,0.003857) (0.00777,0.003890) (0.00779,0.003922) (0.00781,0.003950) (0.00783,0.003974) (0.00786,0.003990) (0.00788,0.003998) (0.00790,0.004000) (0.00792,0.003995) (0.00794,0.003986) (0.00797,0.003975) (0.00799,0.003965) (0.00801,0.003958) (0.00803,0.003957) (0.00805,0.003962) (0.00807,0.003974) (0.00810,0.003993) (0.00812,0.004018) (0.00814,0.004048) (0.00816,0.004080) (0.00818,0.004113) (0.00821,0.004143) (0.00823,0.004169) (0.00825,0.004190) (0.00827,0.004204) (0.00829,0.004211) (0.00831,0.004211) (0.00834,0.004206) (0.00836,0.004197) (0.00838,0.004187) (0.00840,0.004176) (0.00842,0.004168) (0.00844,0.004164) (0.00847,0.004165) (0.00849,0.004172) (0.00851,0.004186) (0.00853,0.004206) (0.00855,0.004232) (0.00858,0.004261) (0.00860,0.004293) (0.00862,0.004326) (0.00864,0.004357) (0.00866,0.004386) (0.00868,0.004409) (0.00871,0.004428) (0.00873,0.004440) (0.00875,0.004446) (0.00877,0.004446) (0.00879,0.004442) (0.00882,0.004434) (0.00884,0.004423) (0.00886,0.004412) (0.00888,0.004403) (0.00890,0.004396) (0.00892,0.004393) (0.00895,0.004396) (0.00897,0.004404) (0.00899,0.004418) (0.00901,0.004437) (0.00903,0.004462) (0.00906,0.004490) (0.00908,0.004521) (0.00910,0.004554) (0.00912,0.004586) (0.00914,0.004617) (0.00916,0.004645) (0.00919,0.004668) (0.00921,0.004687) (0.00923,0.004700) (0.00925,0.004708) (0.00927,0.004710) (0.00930,0.004708) (0.00932,0.004701) (0.00934,0.004692) (0.00936,0.004682) (0.00938,0.004671) (0.00940,0.004661) (0.00943,0.004654) (0.00945,0.004650) (0.00947,0.004651) (0.00949,0.004656) (0.00951,0.004667) (0.00954,0.004683) (0.00956,0.004703) (0.00958,0.004728) (0.00960,0.004757) (0.00962,0.004788) (0.00964,0.004820) (0.00967,0.004853) (0.00969,0.004885) (0.00971,0.004914) (0.00973,0.004940) (0.00975,0.004963) (0.00977,0.004981) (0.00980,0.004994) (0.00982,0.005003) (0.00984,0.005006) (0.00986,0.005005) (0.00988,0.005001) (0.00991,0.004993) (0.00993,0.004983) (0.00995,0.004973) (0.00997,0.004962) (0.00999,0.004952) (0.01001,0.004944) (0.01004,0.004940) (0.01006,0.004939) (0.01008,0.004941) (0.01010,0.004949) (0.01012,0.004961) (0.01015,0.004978) (0.01017,0.004999) (0.01019,0.005023) (0.01021,0.005051) (0.01023,0.005082) (0.01025,0.005114) (0.01028,0.005146) (0.01030,0.005179) (0.01032,0.005210) (0.01034,0.005239) (0.01036,0.005266) (0.01039,0.005289) (0.01041,0.005308) (0.01043,0.005323) (0.01045,0.005334) (0.01047,0.005340) (0.01049,0.005342) (0.01052,0.005341) (0.01054,0.005336) (0.01056,0.005329) (0.01058,0.005319) (0.01060,0.005309) (0.01063,0.005298) (0.01065,0.005287) (0.01067,0.005278) (0.01069,0.005271) (0.01071,0.005266) (0.01073,0.005265) (0.01076,0.005267) (0.01078,0.005273) (0.01080,0.005284) (0.01082,0.005298) (0.01084,0.005316) (0.01086,0.005338) (0.01089,0.005363) (0.01091,0.005391) (0.01093,0.005421) (0.01095,0.005452) (0.01097,0.005485) (0.01100,0.005517) (0.01102,0.005550) (0.01104,0.005580) (0.01106,0.005609) (0.01108,0.005636) (0.01110,0.005660) (0.01113,0.005680) (0.01115,0.005697) (0.01117,0.005710) (0.01119,0.005719) (0.01121,0.005725) (0.01124,0.005727) (0.01126,0.005725) (0.01128,0.005721) (0.01130,0.005714) (0.01132,0.005706) (0.01134,0.005696) (0.01137,0.005685) (0.01139,0.005674) (0.01141,0.005664) (0.01143,0.005654) (0.01145,0.005647) (0.01148,0.005641) (0.01150,0.005638) (0.01152,0.005638) (0.01154,0.005642) (0.01156,0.005648) (0.01158,0.005658) (0.01161,0.005672) (0.01163,0.005689) (0.01165,0.005710) (0.01167,0.005733) (0.01169,0.005759) (0.01172,0.005788) (0.01174,0.005818) (0.01176,0.005849) (0.01178,0.005882) (0.01180,0.005914) (0.01182,0.005947) (0.01185,0.005979) (0.01187,0.006009) (0.01189,0.006038) (0.01191,0.006064) (0.01193,0.006088) (0.01196,0.006109) (0.01198,0.006128) (0.01200,0.006143) (0.01202,0.006154) (0.01204,0.006163) (0.01206,0.006168) (0.01209,0.006170) (0.01211,0.006170) (0.01213,0.006166) (0.01215,0.006161) (0.01217,0.006153) (0.01219,0.006144) (0.01222,0.006134) (0.01224,0.006123) (0.01226,0.006113) (0.01228,0.006102) (0.01230,0.006092) (0.01233,0.006083) (0.01235,0.006076) (0.01237,0.006071) (0.01239,0.006068) (0.01241,0.006068) (0.01243,0.006070) (0.01246,0.006075) (0.01248,0.006084) (0.01250,0.006095) (0.01252,0.006109) (0.01254,0.006126) (0.01257,0.006146) (0.01259,0.006169) (0.01261,0.006194) (0.01263,0.006221) (0.01265,0.006250) (0.01267,0.006281) (0.01270,0.006312) (0.01272,0.006345) (0.01274,0.006378) (0.01276,0.006410) (0.01278,0.006442) (0.01281,0.006474) (0.01283,0.006504) (0.01285,0.006532) (0.01287,0.006559) (0.01289,0.006584) (0.01291,0.006606) (0.01294,0.006626) (0.01296,0.006643) (0.01298,0.006658) (0.01300,0.006670) (0.01302,0.006678) (0.01305,0.006685) (0.01307,0.006688) (0.01309,0.006689) (0.01311,0.006687) (0.01313,0.006684) (0.01315,0.006678) (0.01318,0.006671) (0.01320,0.006662) (0.01322,0.006653) (0.01324,0.006642) (0.01326,0.006631) (0.01328,0.006620) (0.01331,0.006610) (0.01333,0.006600) (0.01335,0.006591) (0.01337,0.006583) (0.01339,0.006576) (0.01342,0.006572) (0.01344,0.006569) (0.01346,0.006568) (0.01348,0.006570) (0.01350,0.006574) (0.01352,0.006581) (0.01355,0.006590) (0.01357,0.006602) (0.01359,0.006616) (0.01361,0.006633) (0.01363,0.006652) (0.01366,0.006674) (0.01368,0.006697) (0.01370,0.006723) (0.01372,0.006750) (0.01374,0.006779) (0.01376,0.006810) (0.01379,0.006841) (0.01381,0.006873) (0.01383,0.006906) (0.01385,0.006938) (0.01387,0.006971) (0.01390,0.007003) (0.01392,0.007035) (0.01394,0.007065) (0.01396,0.007095) (0.01398,0.007123) (0.01400,0.007150) (0.01403,0.007174) (0.01405,0.007197) (0.01407,0.007218) (0.01409,0.007237) (0.01411,0.007253) (0.01414,0.007267) (0.01416,0.007279) (0.01418,0.007288) (0.01420,0.007295) (0.01422,0.007299) (0.01424,0.007302) (0.01427,0.007302) (0.01429,0.007301) (0.01431,0.007297) (0.01433,0.007292) (0.01435,0.007285) (0.01438,0.007277) (0.01440,0.007268) (0.01442,0.007259) (0.01444,0.007248) (0.01446,0.007237) (0.01448,0.007227) (0.01451,0.007216) (0.01453,0.007205) (0.01455,0.007195) (0.01457,0.007186) (0.01459,0.007178) (0.01461,0.007171) (0.01464,0.007166) (0.01466,0.007162) (0.01468,0.007159) (0.01470,0.007159) (0.01472,0.007160) (0.01475,0.007163) (0.01477,0.007169) (0.01479,0.007177) (0.01481,0.007186) (0.01483,0.007198) (0.01485,0.007212) (0.01488,0.007228) (0.01490,0.007247) (0.01492,0.007267) (0.01494,0.007289) (0.01496,0.007313) (0.01499,0.007338) (0.01501,0.007365) (0.01503,0.007394) (0.01505,0.007423) (0.01507,0.007454) (0.01509,0.007485) (0.01512,0.007517) (0.01514,0.007550) (0.01516,0.007582) (0.01518,0.007615) (0.01520,0.007648) (0.01523,0.007680) (0.01525,0.007711) (0.01527,0.007742) (0.01529,0.007772) (0.01531,0.007801) (0.01533,0.007829) (0.01536,0.007856) (0.01538,0.007881) (0.01540,0.007904) (0.01542,0.007926) (0.01544,0.007946) (0.01547,0.007964) (0.01549,0.007980) (0.01551,0.007994) (0.01553,0.008006) (0.01555,0.008017) (0.01557,0.008025) (0.01560,0.008032) (0.01562,0.008036) (0.01564,0.008039) (0.01566,0.008040) (0.01568,0.008039) (0.01571,0.008037) (0.01573,0.008033) (0.01575,0.008028) (0.01577,0.008021) (0.01579,0.008014) (0.01581,0.008005) (0.01584,0.007996) (0.01586,0.007986) (0.01588,0.007976) (0.01590,0.007965) (0.01592,0.007954) (0.01594,0.007943) (0.01597,0.007932) (0.01599,0.007922) (0.01601,0.007912) (0.01603,0.007903) (0.01605,0.007894) (0.01608,0.007886) (0.01610,0.007880) (0.01612,0.007874) (0.01614,0.007870) (0.01616,0.007867) (0.01618,0.007866) (0.01621,0.007866) (0.01623,0.007868) (0.01625,0.007871) (0.01627,0.007876) (0.01629,0.007883) (0.01632,0.007892) (0.01634,0.007902) (0.01636,0.007914) (0.01638,0.007928) (0.01640,0.007944) (0.01642,0.007961) (0.01645,0.007980) (0.01647,0.008001) (0.01649,0.008023) (0.01651,0.008047) (0.01653,0.008072) (0.01656,0.008098) (0.01658,0.008126) (0.01660,0.008154) (0.01662,0.008184) (0.01664,0.008214) (0.01666,0.008245) (0.01669,0.008277) (0.01671,0.008309) (0.01673,0.008341) (0.01675,0.008374) (0.01677,0.008407) (0.01680,0.008439) (0.01682,0.008472) (0.01684,0.008504) (0.01686,0.008536) (0.01688,0.008567) (0.01690,0.008597) (0.01693,0.008627) (0.01695,0.008656) (0.01697,0.008683) (0.01699,0.008710) (0.01701,0.008735) (0.01703,0.008760) (0.01706,0.008782) (0.01708,0.008804) (0.01710,0.008824) (0.01712,0.008842) (0.01714,0.008859) (0.01717,0.008875) (0.01719,0.008889) (0.01721,0.008901) (0.01723,0.008911) (0.01725,0.008920) (0.01727,0.008928) (0.01730,0.008934) (0.01732,0.008938) (0.01734,0.008941) (0.01736,0.008942) (0.01738,0.008942) (0.01741,0.008941) (0.01743,0.008938) (0.01745,0.008934) (0.01747,0.008930) (0.01749,0.008924) (0.01751,0.008917) (0.01754,0.008909) (0.01756,0.008901) (0.01758,0.008891) (0.01760,0.008882) (0.01762,0.008872) (0.01765,0.008861) (0.01767,0.008850) (0.01769,0.008839) (0.01771,0.008828) (0.01773,0.008818) (0.01775,0.008807) (0.01778,0.008797) (0.01780,0.008787) (0.01782,0.008777) (0.01784,0.008768) (0.01786,0.008760) (0.01789,0.008752) (0.01791,0.008746) (0.01793,0.008740) (0.01795,0.008735) (0.01797,0.008732) (0.01799,0.008729) (0.01802,0.008728) (0.01804,0.008728) (0.01806,0.008729) (0.01808,0.008731) (0.01810,0.008735) (0.01813,0.008740) (0.01815,0.008747) (0.01817,0.008755) (0.01819,0.008764) (0.01821,0.008775) (0.01823,0.008788) (0.01826,0.008801) (0.01828,0.008817) (0.01830,0.008833) (0.01832,0.008851) (0.01834,0.008870) (0.01836,0.008891) (0.01839,0.008913) (0.01841,0.008935) (0.01843,0.008960) (0.01845,0.008985) (0.01847,0.009011) (0.01850,0.009038) (0.01852,0.009066) (0.01854,0.009095) (0.01856,0.009125) (0.01858,0.009155) (0.01860,0.009186) (0.01863,0.009217) (0.01865,0.009249) (0.01867,0.009281) (0.01869,0.009314) (0.01871,0.009346) (0.01874,0.009379) (0.01876,0.009412) (0.01878,0.009444) (0.01880,0.009477) (0.01882,0.009509) (0.01884,0.009541) (0.01887,0.009573) (0.01889,0.009604) (0.01891,0.009634) (0.01893,0.009664) (0.01895,0.009693) (0.01898,0.009722) (0.01900,0.009749) (0.01902,0.009776) (0.01904,0.009802) (0.01906,0.009826) (0.01908,0.009850) (0.01911,0.009873) (0.01913,0.009894) (0.01915,0.009915) (0.01917,0.009934) (0.01919,0.009952) (0.01922,0.009969) (0.01924,0.009985) (0.01926,0.009999) (0.01928,0.010012) (0.01930,0.010024) (0.01932,0.010034) (0.01935,0.010043) (0.01937,0.010051) (0.01939,0.010058) (0.01941,0.010063) (0.01943,0.010068) (0.01945,0.010071) (0.01948,0.010073) (0.01950,0.010073) (0.01952,0.010073) (0.01954,0.010072) (0.01956,0.010069) (0.01959,0.010066) (0.01961,0.010062) (0.01963,0.010057) (0.01965,0.010051) (0.01967,0.010044) (0.01969,0.010037) (0.01972,0.010029) (0.01974,0.010020) (0.01976,0.010011) (0.01978,0.010001) (0.01980,0.009991) (0.01983,0.009981) (0.01985,0.009970) (0.01987,0.009960) (0.01989,0.009949) (0.01991,0.009938) (0.01993,0.009927) (0.01996,0.009916) (0.01998,0.009906) (0.02000,0.009895) (0.02000,0.009895) (0.02027,0.009806) (0.02054,0.009846) (0.02080,0.010050) (0.02107,0.010389) (0.02134,0.010785) (0.02161,0.011150) (0.02187,0.011408) (0.02214,0.011524) (0.02241,0.011506) (0.02268,0.011398) (0.02294,0.011268) (0.02321,0.011184) (0.02348,0.011197) (0.02375,0.011334) (0.02401,0.011592) (0.02428,0.011942) (0.02455,0.012337) (0.02482,0.012727) (0.02508,0.013062) (0.02535,0.013308) (0.02562,0.013447) (0.02589,0.013481) (0.02615,0.013426) (0.02642,0.013314) (0.02669,0.013182) (0.02696,0.013067) (0.02722,0.013002) (0.02749,0.013013) (0.02776,0.013113) (0.02803,0.013305) (0.02829,0.013581) (0.02856,0.013923) (0.02883,0.014308) (0.02910,0.014708) (0.02936,0.015098) (0.02963,0.015452) (0.02990,0.015752) (0.03017,0.015982) (0.03043,0.016136) (0.03070,0.016213) (0.03097,0.016220) (0.03124,0.016167) (0.03151,0.016069) (0.03177,0.015945) (0.03204,0.015811) (0.03231,0.015687) (0.03258,0.015590) (0.03284,0.015533) (0.03311,0.015528) (0.03338,0.015582) (0.03365,0.015699) (0.03391,0.015879) (0.03418,0.016119) (0.03445,0.016411) (0.03472,0.016747) (0.03498,0.017117) (0.03525,0.017509) (0.03552,0.017910) (0.03579,0.018308) (0.03605,0.018693) (0.03632,0.019052) (0.03659,0.019378) (0.03686,0.019663) (0.03712,0.019903) (0.03739,0.020093) (0.03766,0.020232) (0.03793,0.020321) (0.03819,0.020362) (0.03846,0.020359) (0.03873,0.020317) (0.03900,0.020241) (0.03926,0.020140) (0.03953,0.020020) (0.03980,0.019889) (0.04007,0.019755) (0.04033,0.019625) (0.04060,0.019506) (0.04087,0.019405) (0.04114,0.019327) (0.04140,0.019278) (0.04167,0.019261) (0.04194,0.019280) (0.04221,0.019337) (0.04247,0.019434) (0.04274,0.019569) (0.04301,0.019744) (0.04328,0.019957) (0.04355,0.020205) (0.04381,0.020486) (0.04408,0.020796) (0.04435,0.021133) (0.04462,0.021491) (0.04488,0.021867) (0.04515,0.022257) (0.04542,0.022654) (0.04569,0.023056) (0.04595,0.023458) (0.04622,0.023855) (0.04649,0.024243) (0.04676,0.024619) (0.04702,0.024978) (0.04729,0.025319) (0.04756,0.025638) (0.04783,0.025933) (0.04809,0.026201) (0.04836,0.026442) (0.04863,0.026655) (0.04890,0.026838) (0.04916,0.026991) (0.04943,0.027115) (0.04970,0.027210) (0.04997,0.027276) (0.05023,0.027315) (0.05050,0.027328) (0.05077,0.027317) (0.05104,0.027283) (0.05130,0.027228) (0.05157,0.027155) (0.05184,0.027066) (0.05211,0.026963) (0.05237,0.026849) (0.05264,0.026726) (0.05291,0.026596) (0.05318,0.026463) (0.05344,0.026328) (0.05371,0.026193) (0.05398,0.026062) (0.05425,0.025936) (0.05452,0.025817) (0.05478,0.025708) (0.05505,0.025609) (0.05532,0.025524) (0.05559,0.025453) (0.05585,0.025397) (0.05612,0.025359) (0.05639,0.025339) (0.05666,0.025337) (0.05692,0.025355) (0.05719,0.025394) (0.05746,0.025453) (0.05773,0.025534) (0.05799,0.025636) (0.05826,0.025758) (0.05853,0.025902) (0.05880,0.026067) (0.05906,0.026252) (0.05933,0.026457) (0.05960,0.026682) (0.05987,0.026925) (0.06013,0.027186) (0.06040,0.027465) (0.06067,0.027759) (0.06094,0.028069) (0.06120,0.028393) (0.06147,0.028729) (0.06174,0.029078) (0.06201,0.029438) (0.06227,0.029807) (0.06254,0.030185) (0.06281,0.030570) (0.06308,0.030961) (0.06334,0.031357) (0.06361,0.031756) (0.06388,0.032158) (0.06415,0.032561) (0.06441,0.032965) (0.06468,0.033367) (0.06495,0.033768) (0.06522,0.034165) (0.06548,0.034558) (0.06575,0.034947) (0.06602,0.035329) (0.06629,0.035705) (0.06656,0.036073) (0.06682,0.036432) (0.06709,0.036782) (0.06736,0.037123) (0.06763,0.037453) (0.06789,0.037772) (0.06816,0.038079) (0.06843,0.038374) (0.06870,0.038657) (0.06896,0.038927) (0.06923,0.039184) (0.06950,0.039427) (0.06977,0.039657) (0.07003,0.039872) (0.07030,0.040074) (0.07057,0.040261) (0.07084,0.040435) (0.07110,0.040594) (0.07137,0.040739) (0.07164,0.040870) (0.07191,0.040987) (0.07217,0.041091) (0.07244,0.041180) (0.07271,0.041257) (0.07298,0.041320) (0.07324,0.041370) (0.07351,0.041408) (0.07378,0.041433) (0.07405,0.041447) (0.07431,0.041449) (0.07458,0.041439) (0.07485,0.041419) (0.07512,0.041389) (0.07538,0.041349) (0.07565,0.041299) (0.07592,0.041240) (0.07619,0.041173) (0.07645,0.041098) (0.07672,0.041015) (0.07699,0.040925) (0.07726,0.040828) (0.07753,0.040725) (0.07779,0.040617) (0.07806,0.040504) (0.07833,0.040386) (0.07860,0.040264) (0.07886,0.040138) (0.07913,0.040009) (0.07940,0.039878) (0.07967,0.039745) (0.07993,0.039610) (0.08020,0.039473) (0.08047,0.039336) (0.08074,0.039199) (0.08100,0.039062) (0.08127,0.038926) (0.08154,0.038790) (0.08181,0.038657) (0.08207,0.038525) (0.08234,0.038395) (0.08261,0.038268) (0.08288,0.038144) (0.08314,0.038024) (0.08341,0.037907) (0.08368,0.037794) (0.08395,0.037686) (0.08421,0.037582) (0.08448,0.037484) (0.08475,0.037391) (0.08502,0.037303) (0.08528,0.037222) (0.08555,0.037146) (0.08582,0.037077) (0.08609,0.037014) (0.08635,0.036959) (0.08662,0.036910) (0.08689,0.036868) (0.08716,0.036833) (0.08742,0.036807) (0.08769,0.036787) (0.08796,0.036776) (0.08823,0.036772) (0.08849,0.036776) (0.08876,0.036789) (0.08903,0.036809) (0.08930,0.036838) (0.08957,0.036876) (0.08983,0.036921) (0.09010,0.036975) (0.09037,0.037038) (0.09064,0.037109) (0.09090,0.037188) (0.09117,0.037276) (0.09144,0.037373) (0.09171,0.037478) (0.09197,0.037591) (0.09224,0.037713) (0.09251,0.037844) (0.09278,0.037982) (0.09304,0.038129) (0.09331,0.038285) (0.09358,0.038448) (0.09385,0.038620) (0.09411,0.038799) (0.09438,0.038987) (0.09465,0.039182) (0.09492,0.039385) (0.09518,0.039596) (0.09545,0.039815) (0.09572,0.040041) (0.09599,0.040274) (0.09625,0.040514) (0.09652,0.040762) (0.09679,0.041016) (0.09706,0.041277) (0.09732,0.041545) (0.09759,0.041820) (0.09786,0.042101) (0.09813,0.042388) (0.09839,0.042681) (0.09866,0.042980) (0.09893,0.043285) (0.09920,0.043596) (0.09946,0.043912) (0.09973,0.044233) (0.10000,0.044560)};
    \addplot[gray,densely dashed,domain=0:0.1] {x/2};
    \addplot[gray,densely dotted,domain=0:0.1] {x/2+x^2};
    \addplot[gray,densely dotted,domain=0:0.1] {x/2-x^2};
    \node[couleurN,font=\small,anchor=west] at (axis cs:0.073,0.022) {$y = f(x)$};
  \end{axis}
\end{tikzpicture}
\end{center}

\noindent Le graphe de $f$, en bleu, oscille entre les paraboles $y = \dfrac{x}{2} \pm x^{2}$, en pointillés, autour de sa tangente en $0$, en tirets. Il a pour tangente une droite de pente $\dfrac{1}{2}$ à l'origine, mais il redescend sur une infinité de petits intervalles qui s'accumulent en $0$ : on y voit, de plus en plus serrés, des allers-retours qui interdisent l'injectivité près de $0$.
%===============================================================================
\section[Fonctions implicites]{Le théorème des fonctions implicites}
%===============================================================================

\subsection{Le cercle}

Le cercle unité est l'ensemble des zéros de $f(x, y) = x^{2} + y^{2} - 1$. Il n'est pas le graphe d'une fonction de $x$, puisque chaque $x \in \left]-1, 1\right[$ a deux ordonnées. Mais \emph{près d'un point} $(a, b)$ du cercle avec $b > 0$, il coïncide avec le graphe de $\varphi(x) = \sqrt{1 - x^{2}}$ : dans un petit rectangle autour de $(a, b)$, l'équation $f(x, y) = 0$ équivaut à $y = \varphi(x)$. Ce qui fait marcher l'argument est la condition $\dpd{f}{y}(a, b) = 2b \neq 0$ : aux points $(\pm 1, 0)$, où elle tombe en défaut, la tangente est verticale, et aucun rectangle ne convient.

% ---- cercle unite ; a = (0.6, 0.8) ; V = ]0.35, 0.8[, W = ]0.5, 1.05[ ;
% ---- phi(0.35) = 0.9367, phi(0.8) = 0.6 : le graphe reste dans W ; -0.8 est hors de W (calcules)
\begin{center}
\begin{tikzpicture}[scale=2.6]
  \draw[gray,->] (-1.35,0) -- (1.45,0) node[below,font=\small] {$x$};
  \draw[gray,->] (0,-1.2) -- (0,1.25) node[left,font=\small] {$y$};
  \fill[couleurQ!12] (0.35,0.5) rectangle (0.8,1.05);
  \draw[couleurQ,thick] (0.35,0.5) rectangle (0.8,1.05);
  \draw[couleurQ,very thick] (0.35,0) -- (0.8,0);
  \draw[couleurQ,very thick] (0,0.5) -- (0,1.05);
  \node[couleurQ,font=\small,below] at (0.575,-0.02) {$V$};
  \node[couleurQ,font=\small,left] at (-0.02,0.775) {$W$};
  \draw[couleurQ,densely dotted] (0.35,0) -- (0.35,0.5) (0.8,0) -- (0.8,0.5) (0,0.5) -- (0.35,0.5) (0,1.05) -- (0.35,1.05);
  \draw[couleurN,thick] (0,0) circle (1);
  \draw[couleurN,line width=2.2pt] (0.35,0.9367) arc[start angle=69.51,end angle=36.87,radius=1];
  \draw[couleurD,thick] (0.3,1.025) -- (0.9,0.575);
  \fill[rougecours] (0.6,0.8) circle (0.9pt) node[above right=-1pt,font=\small] {$(a, b)$};
  \fill[gray] (0.6,-0.8) circle (0.9pt) node[below right=-1pt,font=\small] {$(a, -b)$};
  \fill[rougecours] (1,0) circle (0.9pt);
  \draw[rougecours,densely dashed] (1,-0.45) -- (1,0.45);
  \node[rougecours,font=\small,anchor=west] at (1.03,0.25) {$(1, 0)$ : tangente verticale};
  \node[couleurN,font=\small,anchor=north east] at (-0.72,-0.72) {$x^{2} + y^{2} = 1$};
\end{tikzpicture}
\end{center}

\noindent Dans le rectangle $V \times W$, en violet, le cercle se réduit à l'arc épaissi, graphe de $\varphi(x) = \sqrt{1 - x^{2}}$ au-dessus de $V$ ; le point $(a, -b)$, qui a la même abscisse, est exclu parce que $-b \notin W$. La tangente en orange a pour pente $\varphi'(a) = -\dfrac{a}{b} = -\dfrac{3}{4}$. Au point $(1, 0)$, en rouge, tout rectangle centré contient deux points du cercle de même abscisse $x < 1$, et aucun d'abscisse $x > 1$ : le cercle n'y est pas le graphe d'une fonction de $x$, mais il l'est d'une fonction de $y$.

\subsection{L'énoncé}

Dans $E \times F$, on note $\dd_{1} f$ et $\dd_{2} f$ les différentielles partielles de $f$ par rapport à la première et à la seconde variable (chapitre 2, section 2.4), et l'on munit $E \times F$ de la norme produit, qui en fait un espace de Banach si $E$ et $F$ le sont.

\begin{center}
\begin{tabularx}{15cm}{||X||}
\hline \\
\textbf{\underline{\textcolor{rougecours}{Théorème 2 : Fonctions implicites :}}} \\
Soit $E$, $F$, $G$ trois espaces de Banach, $U$ un ouvert de $E \times F$, $f : U \to G$ de classe $\Cc^{1}$, et $(a, b) \in U$ tel que
\[ f(a, b) = 0 \qquad \text{et} \qquad \dd_{2} f(a, b) \in \Isom(F, G). \]
Il existe un ouvert $V$ de $E$ contenant $a$, un ouvert $W$ de $F$ contenant $b$, avec $V \times W \subset U$, et une application $\varphi : V \to W$ de classe $\Cc^{1}$ tels que
\[ \rmath{\text{pour tout } (x, y) \in V \times W, \qquad f(x, y) = 0 \iff y = \varphi(x)}. \]
De plus, pour tout $x \in V$, $\dd_{2} f(x, \varphi(x)) \in \Isom(F, G)$ et
\[ \rmath{\dd\varphi(x) = -\dd_{2} f\big(x, \varphi(x)\big)^{-1} \circ \dd_{1} f\big(x, \varphi(x)\big)}. \]
\medbreak
\\ \hline
\end{tabularx}
\end{center}

\begin{demo}
\etape{Une application auxiliaire.} On pose $\Phi(x, y) = (x, f(x, y))$, de $U$ dans $E \times G$. Elle est de classe $\Cc^{1}$, et
\[ \dd\Phi(a, b)\cdot(h, k) = \big(h, \dd_{1} f(a, b)\cdot h + \dd_{2} f(a, b)\cdot k\big). \]
Avec $A = \dd_{1} f(a, b)$ et $B = \dd_{2} f(a, b)$, l'équation $\dd\Phi(a, b)\cdot(h, k) = (u, v)$ a pour unique solution $h = u$, $k = B^{-1}\cdot(v - A\cdot u)$, qui dépend continûment de $(u, v)$ : $\dd\Phi(a, b) \in \Isom(E \times F, E \times G)$.

\etape{Inversion locale.} Par le théorème 1, $\Phi$ est un $\Cc^{1}$-difféomorphisme d'un ouvert $\Omega \ni (a, b)$ sur un ouvert $\Omega' \ni (a, 0)$. Les produits de boules ouvertes sont ouverts dans $E \times F$, et l'un d'eux, $V_{1} \times W$, contient $(a, b)$ et est inclus dans $\Omega$ ; la restriction de $\Phi$ à $V_{1} \times W$ est un $\Cc^{1}$-difféomorphisme sur l'ouvert $\Phi(V_{1} \times W)$. Sa réciproque est de la forme $(x, z) \mapsto (x, \psi(x, z))$, puisque $\Phi$ ne change pas la première composante, et $\psi$ est de classe $\Cc^{1}$.

\etape{La fonction implicite.} On pose $V = \{x \in V_{1} \mid (x, 0) \in \Phi(V_{1} \times W)\}$, ouvert comme image réciproque d'un ouvert par $x \mapsto (x, 0)$, qui contient $a$, et $\varphi(x) = \psi(x, 0)$ pour $x \in V$. Alors $(x, \varphi(x)) = \Phi^{-1}(x, 0) \in V_{1} \times W$, donc $\varphi(x) \in W$, et $\varphi$ est de classe $\Cc^{1}$. Pour $(x, y) \in V \times W$,
\[ f(x, y) = 0 \iff \Phi(x, y) = (x, 0) \iff (x, y) = \Phi^{-1}(x, 0) = (x, \varphi(x)), \]
où l'équivalence centrale utilise la bijectivité de $\Phi$ sur $V_{1} \times W$.

\etape{Différentielle.} Soit $x \in V$ et $c = (x, \varphi(x))$. La différentielle $\dd\Phi(c)$ est dans $\Isom(E \times F, E \times G)$ (proposition 6), et $\dd\Phi(c)\cdot(h, k) = (h, \dd_{1} f(c)\cdot h + \dd_{2} f(c)\cdot k)$. Si $\dd_{2} f(c)\cdot k = 0$, alors $\dd\Phi(c)\cdot(0, k) = 0$, donc $k = 0$. Pour $v \in G$, l'antécédent $(h, k)$ de $(0, v)$ par $\dd\Phi(c)$ vérifie $h = 0$ et $\dd_{2} f(c)\cdot k = v$. Ainsi $\dd_{2} f(c)$ est bijective, et sa réciproque, qui associe à $v$ la seconde composante de $\dd\Phi(c)^{-1}\cdot(0, v)$, est continue. Enfin, $x \mapsto f(x, \varphi(x))$ est nulle sur $V$ ; la règle de la chaîne donne $\dd_{1} f(x, \varphi(x)) + \dd_{2} f(x, \varphi(x)) \circ \dd\varphi(x) = 0$, d'où la formule.
\end{demo}

\begin{remarques}
  \item La formule de $\dd\varphi$ se retrouve sans l'apprendre : on dérive l'identité $f(x, \varphi(x)) = 0$ par la règle de la chaîne. L'erreur la plus répandue est d'oublier que les dérivées partielles sont évaluées au point $(x, \varphi(x))$, et non en $(a, b)$.
  \item Le théorème ne donne aucune formule pour $\varphi$, seulement son existence, sa régularité et sa différentielle. Pour en calculer des valeurs, on résout numériquement $f(x, y) = 0$ en $y$, à $x$ fixé, par la méthode de Newton de la section 3.2.
  \item Réciproquement, le théorème d'inversion locale se déduit de celui-ci, appliqué à la fonction $(y, x) \mapsto f(x) - y$, dont on cherche les zéros sous la forme $x = \varphi(y)$ : les deux théorèmes sont deux formes d'un même résultat.
\end{remarques}

\subsection{Le cas de \texorpdfstring{$\R^{p}$}{R\textasciicircum p}}

Prenons $E = \R^{p}$, $F = G = \R^{n}$, et écrivons $f = (f_{1}, \dots, f_{n})$ : on cherche à exprimer $n$ inconnues $y = (y_{1}, \dots, y_{n})$ en fonction de $p$ paramètres $x = (x_{1}, \dots, x_{p})$, à partir de $n$ équations. Les différentielles partielles ont pour matrices les deux blocs de la jacobienne de $f$ :
\[ \Jac_{x} f = \left(\dpd{f_{i}}{x_{j}}\right)_{\substack{1 \leqslant i \leqslant n \\ 1 \leqslant j \leqslant p}} \in \mathcal{M}_{n, p}(\R), \qquad \Jac_{y} f = \left(\dpd{f_{i}}{y_{j}}\right)_{1 \leqslant i, j \leqslant n} \in \Mn{n}. \]

\begin{cadre}
\propo{8} Soit $U$ un ouvert de $\R^{p} \times \R^{n}$, $f : U \to \R^{n}$ de classe $\Cc^{1}$, et $(a, b) \in U$ tel que $f(a, b) = 0$ et $\det \Jac_{y} f(a, b) \neq 0$. Il existe des ouverts $V \ni a$, $W \ni b$ et $\varphi : V \to W$ de classe $\Cc^{1}$ tels que, sur $V \times W$, $f(x, y) = 0$ équivaille à $y = \varphi(x)$, et
\[ \rmath{\Jac_{\varphi}(x) = -\Jac_{y} f\big(x, \varphi(x)\big)^{-1}\,\Jac_{x} f\big(x, \varphi(x)\big)}. \]
En particulier, pour $p = n = 1$ et $\dpd{f}{y}(a, b) \neq 0$ :
\[ \rmath{\varphi'(x) = -\frac{\dfrac{\partial f}{\partial x}\big(x, \varphi(x)\big)}{\dfrac{\partial f}{\partial y}\big(x, \varphi(x)\big)}}. \]
\end{cadre}

\begin{demo}
C'est le théorème 2, la condition $\dd_{2} f(a, b) \in \Isom(\R^{n}, \R^{n})$ équivalant à l'inversibilité de la matrice $\Jac_{y} f(a, b)$ ; la formule traduit en matrices celle du théorème (chapitre 2, proposition 11).
\end{demo}

\begin{exemples}
  \item \emph{Le cercle.} Pour $f(x, y) = x^{2} + y^{2} - 1$ et $b \neq 0$, on retrouve $\varphi'(x) = -\dfrac{2x}{2\varphi(x)} = -\dfrac{x}{\varphi(x)}$, ce que donne aussi la dérivation de $\sqrt{1 - x^{2}}$.
  \item \emph{Une courbe de $\R^{3}$.} L'intersection de la sphère $x^{2} + y^{2} + z^{2} = 3$ et du plan $x + y - z = 1$ passe par $(1, 1, 1)$. Avec $f(x, y, z) = (x^{2} + y^{2} + z^{2} - 3, x + y - z - 1)$, le bloc relatif à $(y, z)$ en ce point est $\begin{pmatrix} 2 & 2 \\ 1 & -1 \end{pmatrix}$, de déterminant $-4 \neq 0$. Près de $(1, 1, 1)$, la courbe est donc le graphe d'une application $x \mapsto (y(x), z(x))$ de classe $\Cc^{1}$, et
  \[ \begin{pmatrix} y'(1) \\ z'(1) \end{pmatrix} = -\begin{pmatrix} 2 & 2 \\ 1 & -1 \end{pmatrix}^{-1}\begin{pmatrix} 2 \\ 1 \end{pmatrix} = \begin{pmatrix} -1 \\ 0 \end{pmatrix}. \]
  Le vecteur $(1, -1, 0)$ dirige la tangente à la courbe en $(1, 1, 1)$ ; il est orthogonal à la fois à $(1, 1, 1)$, normal à la sphère, et à $(1, 1, -1)$, normal au plan, ce que la section 3.5 expliquera.
\end{exemples}

\subsection{Le folium de Descartes}

On ne sait pas résoudre commodément en $y$ l'équation $x^{3} + y^{3} - 3xy = 0$ de cette courbe. Avec $f(x, y) = x^{3} + y^{3} - 3xy$, on a $\dpd{f}{x} = 3(x^{2} - y)$ et $\dpd{f}{y} = 3(y^{2} - x)$. Les points de la courbe où $\dpd{f}{y}$ s'annule vérifient $x = y^{2}$ et $y^{6} + y^{3} - 3y^{3} = 0$, soit $y = 0$ ou $y^{3} = 2$ : ce sont l'origine et le point $\big(2^{2/3}, 2^{1/3}\big) \approx (1{,}587 ; 1{,}260)$. Partout ailleurs, la courbe est localement le graphe d'une fonction $\varphi$ de classe $\Cc^{1}$. Au point $\left(\dfrac{3}{2}, \dfrac{3}{2}\right)$, par exemple, $\dpd{f}{x} = \dpd{f}{y} = \dfrac{9}{4}$, et $\varphi'\left(\dfrac{3}{2}\right) = -1$.

% ---- folium x = 3t/(1+t^3), y = 3t^2/(1+t^3) ; tangente en (3/2, 3/2) de pente -1 ;
% ---- tangente verticale en (2^(2/3), 2^(1/3)) = (1.5874, 1.2599) ; asymptote x + y + 1 = 0 (calcules)
\begin{center}
\begin{tikzpicture}[scale=1.55]
  \draw[gray,->] (-2.3,0) -- (2.4,0) node[below,font=\small] {$x$};
  \draw[gray,->] (0,-2.3) -- (0,2.35) node[left,font=\small] {$y$};
  \draw[gray,densely dashed] (-2.3,1.3) -- (1.3,-2.3);
  \node[gray,font=\small,anchor=north east] at (-1.5,0.32) {$x + y + 1 = 0$};
  \draw[couleurN,thick] plot[smooth] coordinates {(-2.286,1.369) (-2.144,1.238) (-2.013,1.119) (-1.892,1.011) (-1.780,0.913) (-1.674,0.823) (-1.575,0.741) (-1.481,0.665) (-1.392,0.595) (-1.306,0.531) (-1.224,0.471) (-1.145,0.416) (-1.069,0.366) (-0.995,0.319) (-0.923,0.276) (-0.852,0.237) (-0.783,0.201) (-0.715,0.168) (-0.648,0.138) (-0.581,0.112) (-0.516,0.088) (-0.450,0.067) (-0.386,0.049) (-0.321,0.034) (-0.257,0.022) (-0.192,0.012) (-0.128,0.005) (-0.064,0.001) (0.000,0.000) (0.000,0.000) (0.060,0.001) (0.120,0.005) (0.180,0.011) (0.240,0.019) (0.300,0.030) (0.360,0.043) (0.420,0.059) (0.479,0.077) (0.538,0.097) (0.597,0.120) (0.655,0.144) (0.712,0.171) (0.768,0.200) (0.824,0.231) (0.878,0.264) (0.932,0.299) (0.984,0.335) (1.034,0.373) (1.083,0.413) (1.130,0.453) (1.175,0.495) (1.219,0.538) (1.260,0.581) (1.299,0.625) (1.336,0.669) (1.370,0.714) (1.402,0.759) (1.431,0.803) (1.458,0.848) (1.482,0.891) (1.504,0.935) (1.523,0.977) (1.539,1.018) (1.553,1.059) (1.565,1.098) (1.574,1.136) (1.580,1.172) (1.585,1.207) (1.587,1.241) (1.587,1.273) (1.585,1.303) (1.582,1.332) (1.576,1.359) (1.569,1.384) (1.561,1.408) (1.550,1.430) (1.539,1.450) (1.526,1.469) (1.513,1.486) (1.498,1.502) (1.483,1.516) (1.466,1.529) (1.449,1.540) (1.431,1.550) (1.413,1.558) (1.394,1.566) (1.375,1.572) (1.356,1.577) (1.336,1.581) (1.317,1.584) (1.297,1.586) (1.277,1.587) (1.257,1.587) (1.237,1.587) (1.217,1.586) (1.197,1.584) (1.177,1.581) (1.157,1.578) (1.138,1.574) (1.118,1.570) (1.099,1.565) (1.080,1.560) (1.062,1.554) (1.043,1.548) (1.025,1.542) (1.007,1.535) (0.990,1.528) (0.972,1.521) (0.955,1.513) (0.939,1.506) (0.922,1.498) (0.906,1.490) (0.890,1.481) (0.875,1.473) (0.859,1.464) (0.844,1.456) (0.830,1.447) (0.815,1.438) (0.801,1.430) (0.787,1.421) (0.774,1.412) (0.761,1.403) (0.748,1.394) (0.735,1.385) (0.722,1.376) (0.710,1.367) (0.698,1.358) (0.687,1.349) (0.675,1.340) (0.664,1.331) (0.653,1.322) (0.642,1.313) (0.632,1.305) (0.621,1.296) (0.611,1.287) (0.602,1.278) (0.592,1.270) (0.582,1.261) (0.573,1.253) (0.564,1.244) (0.555,1.236) (0.547,1.228) (0.538,1.219) (0.530,1.211) (0.522,1.203) (0.514,1.195) (0.506,1.187) (0.498,1.179) (0.491,1.171) (0.484,1.163) (0.476,1.156) (0.469,1.148) (0.462,1.140) (0.456,1.133) (0.449,1.126) (0.443,1.118) (0.436,1.111) (0.430,1.104) (0.424,1.097) (0.418,1.089) (0.412,1.082) (0.406,1.076) (0.401,1.069) (0.395,1.062) (0.390,1.055) (0.385,1.048) (0.379,1.042) (0.374,1.035) (0.369,1.029) (0.364,1.023) (0.359,1.016) (0.355,1.010) (0.350,1.004) (0.346,0.998) (0.341,0.992) (0.337,0.986) (0.332,0.980) (0.328,0.974) (0.324,0.968) (0.320,0.962) (0.316,0.956) (0.312,0.951) (0.308,0.945) (0.304,0.940) (0.301,0.934) (0.297,0.929) (0.293,0.923) (0.290,0.918) (0.286,0.913) (0.283,0.908) (0.280,0.903) (0.276,0.897) (0.273,0.892) (0.270,0.887) (0.267,0.882) (0.264,0.878) (0.261,0.873) (0.258,0.868) (0.255,0.863) (0.252,0.858) (0.249,0.854) (0.246,0.849) (0.244,0.845) (0.241,0.840) (0.238,0.836) (0.236,0.831) (0.233,0.827) (0.230,0.822) (0.228,0.818) (0.226,0.814) (0.223,0.810) (0.221,0.806) (0.218,0.801) (0.216,0.797) (0.214,0.793) (0.212,0.789) (0.209,0.785) (0.207,0.781) (0.205,0.777) (0.203,0.774) (0.201,0.770) (0.199,0.766) (0.197,0.762) (0.195,0.758) (0.193,0.755) (0.191,0.751) (0.189,0.747) (0.187,0.744) (0.186,0.740) (0.184,0.737) (0.182,0.733) (0.180,0.730) (0.178,0.726) (0.177,0.723) (0.175,0.720) (0.173,0.716) (0.172,0.713) (0.170,0.710) (0.169,0.706) (0.167,0.703) (0.165,0.700) (0.164,0.697) (0.162,0.694) (0.161,0.690) (0.159,0.687) (0.158,0.684) (0.157,0.681) (0.155,0.678) (0.154,0.675) (0.152,0.672) (0.151,0.669) (0.150,0.666) (0.148,0.664) (0.147,0.661) (0.146,0.658) (0.145,0.655) (0.143,0.652) (0.142,0.649) (0.141,0.647) (0.140,0.644) (0.138,0.641) (0.137,0.639) (0.136,0.636) (0.135,0.633) (0.134,0.631) (0.133,0.628) (0.132,0.625) (0.131,0.623) (0.129,0.620) (0.128,0.618) (0.127,0.615) (0.126,0.613) (0.125,0.610) (0.124,0.608) (0.123,0.606) (0.122,0.603) (0.121,0.601) (0.120,0.598) (0.119,0.596) (0.118,0.594) (0.118,0.591) (0.117,0.589) (0.116,0.587) (0.115,0.585) (0.114,0.582) (0.113,0.580) (0.112,0.578) (0.111,0.576) (0.110,0.574) (0.110,0.571) (0.109,0.569) (0.108,0.567) (0.107,0.565) (0.106,0.563) (0.106,0.561) (0.105,0.559) (0.104,0.557) (0.103,0.555) (0.102,0.553) (0.102,0.551) (0.101,0.549) (0.100,0.547) (0.100,0.545) (0.099,0.543) (0.098,0.541) (0.097,0.539) (0.097,0.537) (0.096,0.535) (0.095,0.533) (0.095,0.531) (0.094,0.530) (0.093,0.528) (0.093,0.526) (0.092,0.524) (0.091,0.522) (0.091,0.520) (0.090,0.519) (0.090,0.517) (0.089,0.515) (0.088,0.513) (0.088,0.512) (0.087,0.510) (0.086,0.508) (0.086,0.506) (0.085,0.505) (0.085,0.503) (0.084,0.501) (0.084,0.500) (0.083,0.498) (0.083,0.496) (0.082,0.495) (0.081,0.493) (0.081,0.492) (0.080,0.490) (0.080,0.488) (0.079,0.487) (0.079,0.485) (0.078,0.484) (0.078,0.482) (0.077,0.481) (0.077,0.479) (0.076,0.478) (0.076,0.476) (0.075,0.475) (0.075,0.473) (0.074,0.472) (0.074,0.470) (0.074,0.469) (0.073,0.467) (0.073,0.466) (0.072,0.464) (0.072,0.463) (0.071,0.462) (0.071,0.460) (0.070,0.459) (0.070,0.457) (0.070,0.456) (0.069,0.455) (0.069,0.453) (0.068,0.452) (0.068,0.450) (0.067,0.449) (0.067,0.448) (0.067,0.446) (0.066,0.445) (0.066,0.444) (0.065,0.443) (0.065,0.441) (0.065,0.440) (0.064,0.439) (0.064,0.437) (0.064,0.436) (0.063,0.435) (0.063,0.434) (0.063,0.432) (0.062,0.431) (0.062,0.430) (0.061,0.429) (0.061,0.427) (0.061,0.426) (0.060,0.425) (0.060,0.424) (0.060,0.423) (0.059,0.421) (0.059,0.420) (0.059,0.419) (0.058,0.418) (0.058,0.417) (0.058,0.416) (0.057,0.415) (0.057,0.413) (0.057,0.412) (0.056,0.411) (0.056,0.410) (0.056,0.409) (0.056,0.408) (0.055,0.407) (0.055,0.406) (0.055,0.404) (0.054,0.403) (0.054,0.402) (0.054,0.401) (0.054,0.400) (0.053,0.399) (0.053,0.398) (0.053,0.397) (0.052,0.396) (0.052,0.395) (0.052,0.394) (0.052,0.393) (0.051,0.392) (0.051,0.391) (0.051,0.390) (0.050,0.389) (0.050,0.388) (0.050,0.387) (0.050,0.386) (0.049,0.385) (0.049,0.384) (0.049,0.383) (0.049,0.382) (0.048,0.381) (0.048,0.380) (0.048,0.379) (0.048,0.378) (0.047,0.377) (0.047,0.376) (0.047,0.375) (0.047,0.374) (0,0)};
  \draw[couleurN,thick] plot[smooth] coordinates {(0,0) (0.021,-0.250) (0.021,-0.251) (0.021,-0.252) (0.021,-0.253) (0.022,-0.255) (0.022,-0.256) (0.022,-0.257) (0.022,-0.258) (0.022,-0.259) (0.023,-0.260) (0.023,-0.262) (0.023,-0.263) (0.023,-0.264) (0.023,-0.265) (0.024,-0.266) (0.024,-0.268) (0.024,-0.269) (0.024,-0.270) (0.025,-0.271) (0.025,-0.273) (0.025,-0.274) (0.025,-0.275) (0.026,-0.277) (0.026,-0.278) (0.026,-0.279) (0.026,-0.281) (0.027,-0.282) (0.027,-0.284) (0.027,-0.285) (0.027,-0.286) (0.028,-0.288) (0.028,-0.289) (0.028,-0.291) (0.028,-0.292) (0.029,-0.294) (0.029,-0.295) (0.029,-0.297) (0.030,-0.298) (0.030,-0.300) (0.030,-0.301) (0.031,-0.303) (0.031,-0.305) (0.031,-0.306) (0.032,-0.308) (0.032,-0.310) (0.032,-0.311) (0.033,-0.313) (0.033,-0.315) (0.033,-0.316) (0.034,-0.318) (0.034,-0.320) (0.034,-0.322) (0.035,-0.324) (0.035,-0.325) (0.036,-0.327) (0.036,-0.329) (0.036,-0.331) (0.037,-0.333) (0.037,-0.335) (0.038,-0.337) (0.038,-0.339) (0.039,-0.341) (0.039,-0.343) (0.040,-0.345) (0.040,-0.347) (0.041,-0.349) (0.041,-0.351) (0.042,-0.354) (0.042,-0.356) (0.043,-0.358) (0.043,-0.360) (0.044,-0.363) (0.044,-0.365) (0.045,-0.367) (0.045,-0.370) (0.046,-0.372) (0.047,-0.374) (0.047,-0.377) (0.048,-0.379) (0.049,-0.382) (0.049,-0.384) (0.050,-0.387) (0.051,-0.390) (0.051,-0.392) (0.052,-0.395) (0.053,-0.398) (0.053,-0.401) (0.054,-0.403) (0.055,-0.406) (0.056,-0.409) (0.056,-0.412) (0.057,-0.415) (0.058,-0.418) (0.059,-0.421) (0.060,-0.424) (0.061,-0.428) (0.062,-0.431) (0.063,-0.434) (0.064,-0.437) (0.065,-0.441) (0.066,-0.444) (0.067,-0.448) (0.068,-0.451) (0.069,-0.455) (0.070,-0.459) (0.071,-0.462) (0.072,-0.466) (0.073,-0.470) (0.075,-0.474) (0.076,-0.478) (0.077,-0.482) (0.078,-0.486) (0.080,-0.490) (0.081,-0.494) (0.083,-0.499) (0.084,-0.503) (0.085,-0.508) (0.087,-0.512) (0.089,-0.517) (0.090,-0.522) (0.092,-0.526) (0.094,-0.531) (0.095,-0.536) (0.097,-0.542) (0.099,-0.547) (0.101,-0.552) (0.103,-0.558) (0.105,-0.563) (0.107,-0.569) (0.109,-0.574) (0.111,-0.580) (0.114,-0.586) (0.116,-0.593) (0.119,-0.599) (0.121,-0.605) (0.124,-0.612) (0.126,-0.618) (0.129,-0.625) (0.132,-0.632) (0.135,-0.640) (0.138,-0.647) (0.141,-0.654) (0.145,-0.662) (0.148,-0.670) (0.152,-0.678) (0.155,-0.686) (0.159,-0.695) (0.163,-0.704) (0.167,-0.713) (0.171,-0.722) (0.176,-0.731) (0.180,-0.741) (0.185,-0.751) (0.190,-0.761) (0.195,-0.772) (0.201,-0.783) (0.206,-0.794) (0.212,-0.805) (0.218,-0.817) (0.225,-0.829) (0.231,-0.842) (0.238,-0.855) (0.246,-0.869) (0.254,-0.883) (0.262,-0.897) (0.270,-0.912) (0.279,-0.928) (0.288,-0.944) (0.298,-0.960) (0.309,-0.978) (0.320,-0.996) (0.331,-1.015) (0.343,-1.034) (0.356,-1.055) (0.370,-1.076) (0.385,-1.098) (0.400,-1.122) (0.417,-1.146) (0.434,-1.172) (0.453,-1.199) (0.473,-1.227) (0.495,-1.257) (0.518,-1.289) (0.543,-1.323) (0.570,-1.359) (0.599,-1.397) (0.631,-1.438) (0.665,-1.481) (0.703,-1.528) (0.744,-1.578) (0.789,-1.633) (0.839,-1.692) (0.894,-1.757) (0.956,-1.829) (1.025,-1.908) (1.103,-1.995) (1.192,-2.094) (1.294,-2.205)};
  \draw[couleurD,thick] (0.95,2.05) -- (2.05,0.95);
  \fill[rougecours] (1.5,1.5) circle (1.1pt) node[above right=0pt,font=\small] {$\left(\dfrac{3}{2}, \dfrac{3}{2}\right)$};
  \draw[rougecours,densely dashed] (1.5874,0.75) -- (1.5874,1.75);
  \fill[rougecours] (1.5874,1.2599) circle (1.1pt);
  \node[rougecours,font=\small,anchor=west] at (1.65,0.85) {$\big(2^{2/3}, 2^{1/3}\big)$};
  \fill[rougecours] (0,0) circle (1.1pt);
  \node[rougecours,font=\small,anchor=north west] at (0.05,-0.05) {$O$};
\end{tikzpicture}
\end{center}

\noindent Le folium, en bleu, est localement un graphe au-dessus de l'axe des abscisses, sauf en deux points rouges : au point $\big(2^{2/3}, 2^{1/3}\big)$, la tangente est verticale, en pointillés, et la boucle se replie ; à l'origine, la courbe se recoupe, et aucun rectangle ne l'isole en un seul graphe. En $\left(\dfrac{3}{2}, \dfrac{3}{2}\right)$, la tangente, en orange, a pour pente $-1$, comme le donne la proposition 8.

\begin{remarques}
  \item À l'origine, les deux dérivées partielles sont nulles : ni $x$ ni $y$ ne peuvent servir de paramètre. Au point $\big(2^{2/3}, 2^{1/3}\big)$, seule $\dpd{f}{y}$ s'annule : la courbe y est le graphe d'une fonction de $y$, en échangeant les rôles des variables.
  \item La condition $\dpd{f}{y}(a, b) \neq 0$ est suffisante, non nécessaire : la courbe $y^{3} = x$ est le graphe de $x \mapsto x^{1/3}$, bien que $\dpd{f}{y}$, pour $f(x, y) = y^{3} - x$, s'annule en $(0, 0)$. Mais alors $\varphi$ n'est pas dérivable en $0$.
\end{remarques}

\subsection{Tracer une courbe implicite avec Python}

Les logiciels tracent une courbe $f(x, y) = 0$ par \emph{continuation} : on part d'un point $(x_{0}, y_{0})$ de la courbe, on avance $x$ d'un petit pas, et l'on corrige $y$ par la méthode de Newton en la seule variable $y$, en partant de l'ordonnée précédente. Le théorème des fonctions implicites garantit que, loin des points où $\dpd{f}{y}$ s'annule, la solution cherchée existe, reste proche de la précédente, et que la méthode converge si le pas est assez petit.

\begin{lstlisting}
def f(x, y):  return x**3 + y**3 - 3*x*y
def fy(x, y): return 3*y**2 - 3*x          # derivee partielle en y

def branche(x0, y0, pas, n):
    """Suit la courbe f = 0 depuis (x0, y0) en faisant varier x."""
    points, y = [(x0, y0)], y0
    for k in range(1, n + 1):
        x = x0 + k*pas
        for _ in range(20):                # Newton en y, x fixe
            y = y - f(x, y)/fy(x, y)
        points.append((x, y))
    return points

for x, y in branche(1.5, 1.5, -0.1, 5):
    print(f"{x:.1f}  {y:.6f}")
\end{lstlisting}

\noindent Partant de $\left(\dfrac{3}{2}, \dfrac{3}{2}\right)$ vers la gauche, le programme affiche, pour $x = 1{,}4$ ; $1{,}3$ ; $1{,}2$ ; $1{,}1$ ; $1{,}0$, les ordonnées $1{,}563\,715$ ; $1{,}585\,721$ ; $1{,}584\,014$ ; $1{,}565\,114$ ; $1{,}532\,089$, points de l'arc supérieur de la boucle. Comme $\varphi'\left(\dfrac{3}{2}\right) = -1$, les ordonnées commencent par croître quand $x$ diminue ; elles passent par un maximum voisin de $2^{2/3} \approx 1{,}587$ près de $x = 2^{1/3} \approx 1{,}260$, où la tangente est horizontale, puis décroissent. Vers la droite, avec un pas de $0{,}02$, la continuation fonctionne jusqu'à $x = 1{,}58$, puis échoue près de $x = 2^{2/3} \approx 1{,}587$, où la tangente est verticale : en $x = 1{,}60$, après vingt itérations, $f$ vaut encore $0{,}05$, et en $x = 1{,}62$, la méthode converge vers $y \approx -2{,}554$, un point de la branche inférieure. Au-delà de la tangente verticale, la boucle n'a plus de point, et Newton saute d'une branche à l'autre sans prévenir.
\newpage
%===============================================================================
\section[Hyperplan tangent]{Hyperplan tangent à une hypersurface}
%===============================================================================

Dans cette section, $U$ est un ouvert de $\R^{p}$, muni du produit scalaire usuel $\ps{\cdot}{\cdot}$, et $g : U \to \R$ est de classe $\Cc^{1}$. Une seule équation $g(x) = 0$ entre $p$ variables laisse $p - 1$ degrés de liberté : on attend une «~courbe~» dans $\R^{2}$, une «~surface~» dans $\R^{3}$. Le théorème des fonctions implicites justifie cette intuition, à condition que le gradient ne s'annule pas.

\subsection{Hypersurfaces et points réguliers}

\begin{center}
\begin{tabularx}{15cm}{||X}
\noindent \textbf{\underline{Définition :}} L'ensemble $S = \{x \in U \mid g(x) = 0\}$ est l'\mot{hypersurface} d'équation $g = 0$. Un point $a \in S$ est \mot{régulier} si $\rmath{\grad g(a) \neq 0}$, et \mot{singulier} sinon.

\noindent Un vecteur $v \in \R^{p}$ est \mot{tangent à $S$ en $a$} s'il existe $\varepsilon > 0$ et une application $\gamma : \left]-\varepsilon, \varepsilon\right[ \to \R^{p}$ dérivable en $0$, à valeurs dans $S$, telle que $\gamma(0) = a$ et $\gamma'(0) = v$. On note $\rmath{T_{a}S}$ l'ensemble des vecteurs tangents à $S$ en $a$.
\end{tabularx}
\end{center}

\begin{remarques}
  \item Pour $p = 2$, une hypersurface est une courbe plane, et pour $p = 3$, une surface. Les lignes de niveau $\{f = c\}$ d'une fonction $f$ sont des hypersurfaces, d'équation $f - c = 0$.
  \item Le vecteur nul est toujours tangent, grâce à la courbe constante $\gamma(t) = a$, et si $v$ est tangent, $\lambda v$ l'est aussi, grâce à $t \mapsto \gamma(\lambda t)$. Mais rien n'assure a priori que la somme de deux vecteurs tangents soit tangente : c'est ce qu'apporte la régularité.
\end{remarques}

\begin{cadre}
\propo{9} Soit $a \in S$ un point régulier, et $j$ un indice tel que $\dpd{g}{x_{j}}(a) \neq 0$. Au voisinage de $a$, l'hypersurface $S$ est le graphe d'une fonction de classe $\Cc^{1}$ qui exprime $x_{j}$ en fonction des $p - 1$ autres coordonnées.
\end{cadre}

\begin{demo}
Quitte à renuméroter les coordonnées, $j = p$. On écrit $x = (x', x_{p})$ avec $x' \in \R^{p - 1}$, et l'on applique la proposition 8 avec $n = 1$ : la matrice $\Jac_{y} g(a)$ est le réel $\dpd{g}{x_{p}}(a) \neq 0$. Il existe des ouverts $V \ni a'$ et $W \ni a_{p}$, et $\varphi : V \to W$ de classe $\Cc^{1}$, tels que $S \cap (V \times W) = \{(x', \varphi(x')) \mid x' \in V\}$.
\end{demo}

\subsection{Le théorème de l'hyperplan tangent}

\begin{center}
\begin{tabularx}{15cm}{||X||}
\hline \\
\textbf{\underline{\textcolor{rougecours}{Théorème 3 : Hyperplan tangent :}}} \\
Soit $a$ un point régulier de l'hypersurface $S$ d'équation $g = 0$. L'ensemble des vecteurs tangents à $S$ en $a$ est l'hyperplan orthogonal au gradient :
\[ \rmath{T_{a}S = \grad g(a)^{\perp} = \Ker \dd g(a)}. \]
L'\mot{hyperplan tangent} à $S$ en $a$ est l'hyperplan affine $a + T_{a}S$, d'équation
\[ \rmath{\ps{\grad g(a)}{x - a} = \sum_{j = 1}^{p} \dpd{g}{x_{j}}(a)\,(x_{j} - a_{j}) = 0}. \]
\medbreak
\\ \hline
\end{tabularx}
\end{center}

\begin{demo}
\etape{Les vecteurs tangents sont orthogonaux au gradient.} Si $\gamma$ est à valeurs dans $S$, dérivable en $0$, avec $\gamma(0) = a$ et $\gamma'(0) = v$, alors $g \circ \gamma$ est nulle, et la règle de la chaîne (chapitre 2, proposition 11) donne $0 = (g \circ \gamma)'(0) = \dd g(a)\cdot v = \ps{\grad g(a)}{v}$.

\etape{Réciproquement.} Par la proposition 9, en renumérotant les coordonnées, $S$ coïncide près de $a$ avec le graphe de $\varphi : V \to W$, où $\dpd{g}{x_{p}}(a) \neq 0$. Soit $v = (v', v_{p}) \in \grad g(a)^{\perp}$. On pose $\gamma(t) = (a' + tv', \varphi(a' + tv'))$, défini pour $t$ petit puisque $V$ est ouvert : $\gamma$ est à valeurs dans $S$, $\gamma(0) = a$, et $\gamma'(0) = (v', \dd\varphi(a')\cdot v')$. La proposition 8 donne
\[ \dd\varphi(a')\cdot v' = -\frac{1}{\dpd{g}{x_{p}}(a)}\sum_{j = 1}^{p - 1} \dpd{g}{x_{j}}(a)\,v_{j} = v_{p}, \]
la dernière égalité traduisant $\ps{\grad g(a)}{v} = 0$. Ainsi $\gamma'(0) = v$, et $v$ est tangent.

\etape{Conclusion.} Comme $\grad g(a) \neq 0$, son orthogonal est un hyperplan, noyau de la forme linéaire non nulle $\dd g(a) = \ps{\grad g(a)}{\cdot}$.
\end{demo}

\begin{remarques}
  \item Le gradient $\grad g(a)$ est un \mot{vecteur normal} à $S$ en $a$. Avec la direction de plus forte pente du chapitre 2 (proposition 10), on obtient l'image d'une montagne dont $S$ est une courbe de niveau : la pente la plus raide est perpendiculaire aux courbes de niveau.
  \item \emph{Piège.} En un point singulier, $T_{a}S$ n'est pas forcément un hyperplan, et l'équation de l'hyperplan tangent, dont tous les coefficients sont alors nuls, est satisfaite par tous les points : elle ne décrit rien. L'erreur la plus répandue est d'écrire l'équation de l'hyperplan tangent sans avoir vérifié que le gradient est non nul.
\end{remarques}

\begin{cadre}
\propo{10} \emph{(Gradient et lignes de niveau.)} Soit $f : U \to \R$ de classe $\Cc^{1}$ et $a \in U$ tel que $\grad f(a) \neq 0$. Le gradient $\grad f(a)$ est orthogonal à la ligne de niveau $\{f = f(a)\}$ en $a$ : il est orthogonal à tout vecteur tangent en $a$ à cette ligne de niveau, et ces vecteurs tangents forment l'hyperplan $\grad f(a)^{\perp}$.
\end{cadre}

\begin{demo}
C'est le théorème 3, appliqué à $g = f - f(a)$, qui a le même gradient que $f$.
\end{demo}

\noindent C'est la propriété annoncée au chapitre 2 et observée sur les ellipses de niveau de $\dfrac{x^{2}}{4} + y^{2}$. La figure suivante montre, au contraire, ce qui se passe en un point singulier.

% ---- lemniscate (x^2+y^2)^2 = 2(x^2-y^2), r^2 = 2 cos(2 theta) ; points reguliers :
% ---- (1.4142, 0) grad (5.657, 0) ; (0.866, 0.5) grad (0, 4) ; (1.0987, -0.4551) grad (1.820, -4.395) unitaire (0.3827, -0.9239) ;
% ---- et symetriques ; origine : grad nul (calcules)
\begin{center}
\begin{tikzpicture}[scale=2.5]
  \draw[gray,->] (-1.85,0) -- (1.9,0) node[below,font=\small] {$x$};
  \draw[gray,->] (0,-0.85) -- (0,0.95) node[left,font=\small] {$y$};
  \draw[rougecours,densely dashed] (-0.85,-0.85) -- (0.85,0.85);
  \draw[rougecours,densely dashed] (-0.85,0.85) -- (0.85,-0.85);
  \draw[couleurN,thick] plot[smooth] coordinates {(0.000,-0.000) (0.126,-0.124) (0.180,-0.174) (0.222,-0.212) (0.258,-0.243) (0.291,-0.269) (0.321,-0.292) (0.349,-0.313) (0.376,-0.331) (0.401,-0.348) (0.425,-0.363) (0.449,-0.377) (0.472,-0.390) (0.494,-0.402) (0.515,-0.413) (0.537,-0.423) (0.557,-0.432) (0.577,-0.441) (0.597,-0.449) (0.617,-0.456) (0.636,-0.462) (0.655,-0.468) (0.673,-0.473) (0.692,-0.478) (0.710,-0.482) (0.727,-0.486) (0.745,-0.489) (0.762,-0.492) (0.779,-0.494) (0.796,-0.496) (0.812,-0.498) (0.829,-0.499) (0.845,-0.500) (0.861,-0.500) (0.876,-0.500) (0.892,-0.499) (0.907,-0.499) (0.922,-0.498) (0.937,-0.496) (0.952,-0.494) (0.966,-0.492) (0.980,-0.490) (0.994,-0.487) (1.008,-0.484) (1.022,-0.481) (1.035,-0.477) (1.048,-0.473) (1.061,-0.469) (1.074,-0.465) (1.086,-0.460) (1.099,-0.455) (1.111,-0.450) (1.123,-0.444) (1.134,-0.439) (1.146,-0.433) (1.157,-0.427) (1.168,-0.421) (1.179,-0.414) (1.189,-0.407) (1.200,-0.400) (1.210,-0.393) (1.220,-0.386) (1.229,-0.378) (1.239,-0.370) (1.248,-0.363) (1.257,-0.354) (1.266,-0.346) (1.274,-0.338) (1.282,-0.329) (1.290,-0.321) (1.298,-0.312) (1.306,-0.303) (1.313,-0.293) (1.320,-0.284) (1.327,-0.275) (1.333,-0.265) (1.340,-0.256) (1.346,-0.246) (1.351,-0.236) (1.357,-0.226) (1.362,-0.216) (1.367,-0.206) (1.372,-0.195) (1.377,-0.185) (1.381,-0.174) (1.385,-0.164) (1.389,-0.153) (1.392,-0.143) (1.395,-0.132) (1.398,-0.121) (1.401,-0.110) (1.404,-0.099) (1.406,-0.088) (1.408,-0.077) (1.410,-0.066) (1.411,-0.055) (1.412,-0.044) (1.413,-0.033) (1.414,-0.022) (1.414,-0.011) (1.414,0.000) (1.414,0.011) (1.414,0.022) (1.413,0.033) (1.412,0.044) (1.411,0.055) (1.410,0.066) (1.408,0.077) (1.406,0.088) (1.404,0.099) (1.401,0.110) (1.398,0.121) (1.395,0.132) (1.392,0.143) (1.389,0.153) (1.385,0.164) (1.381,0.174) (1.377,0.185) (1.372,0.195) (1.367,0.206) (1.362,0.216) (1.357,0.226) (1.351,0.236) (1.346,0.246) (1.340,0.256) (1.333,0.265) (1.327,0.275) (1.320,0.284) (1.313,0.293) (1.306,0.303) (1.298,0.312) (1.290,0.321) (1.282,0.329) (1.274,0.338) (1.266,0.346) (1.257,0.354) (1.248,0.363) (1.239,0.370) (1.229,0.378) (1.220,0.386) (1.210,0.393) (1.200,0.400) (1.189,0.407) (1.179,0.414) (1.168,0.421) (1.157,0.427) (1.146,0.433) (1.134,0.439) (1.123,0.444) (1.111,0.450) (1.099,0.455) (1.086,0.460) (1.074,0.465) (1.061,0.469) (1.048,0.473) (1.035,0.477) (1.022,0.481) (1.008,0.484) (0.994,0.487) (0.980,0.490) (0.966,0.492) (0.952,0.494) (0.937,0.496) (0.922,0.498) (0.907,0.499) (0.892,0.499) (0.876,0.500) (0.861,0.500) (0.845,0.500) (0.829,0.499) (0.812,0.498) (0.796,0.496) (0.779,0.494) (0.762,0.492) (0.745,0.489) (0.727,0.486) (0.710,0.482) (0.692,0.478) (0.673,0.473) (0.655,0.468) (0.636,0.462) (0.617,0.456) (0.597,0.449) (0.577,0.441) (0.557,0.432) (0.537,0.423) (0.515,0.413) (0.494,0.402) (0.472,0.390) (0.449,0.377) (0.425,0.363) (0.401,0.348) (0.376,0.331) (0.349,0.313) (0.321,0.292) (0.291,0.269) (0.258,0.243) (0.222,0.212) (0.180,0.174) (0.126,0.124) (0.000,0.000)};
  \draw[couleurN,thick] plot[smooth] coordinates {(-0.000,0.000) (-0.126,0.124) (-0.180,0.174) (-0.222,0.212) (-0.258,0.243) (-0.291,0.269) (-0.321,0.292) (-0.349,0.313) (-0.376,0.331) (-0.401,0.348) (-0.425,0.363) (-0.449,0.377) (-0.472,0.390) (-0.494,0.402) (-0.515,0.413) (-0.537,0.423) (-0.557,0.432) (-0.577,0.441) (-0.597,0.449) (-0.617,0.456) (-0.636,0.462) (-0.655,0.468) (-0.673,0.473) (-0.692,0.478) (-0.710,0.482) (-0.727,0.486) (-0.745,0.489) (-0.762,0.492) (-0.779,0.494) (-0.796,0.496) (-0.812,0.498) (-0.829,0.499) (-0.845,0.500) (-0.861,0.500) (-0.876,0.500) (-0.892,0.499) (-0.907,0.499) (-0.922,0.498) (-0.937,0.496) (-0.952,0.494) (-0.966,0.492) (-0.980,0.490) (-0.994,0.487) (-1.008,0.484) (-1.022,0.481) (-1.035,0.477) (-1.048,0.473) (-1.061,0.469) (-1.074,0.465) (-1.086,0.460) (-1.099,0.455) (-1.111,0.450) (-1.123,0.444) (-1.134,0.439) (-1.146,0.433) (-1.157,0.427) (-1.168,0.421) (-1.179,0.414) (-1.189,0.407) (-1.200,0.400) (-1.210,0.393) (-1.220,0.386) (-1.229,0.378) (-1.239,0.370) (-1.248,0.363) (-1.257,0.354) (-1.266,0.346) (-1.274,0.338) (-1.282,0.329) (-1.290,0.321) (-1.298,0.312) (-1.306,0.303) (-1.313,0.293) (-1.320,0.284) (-1.327,0.275) (-1.333,0.265) (-1.340,0.256) (-1.346,0.246) (-1.351,0.236) (-1.357,0.226) (-1.362,0.216) (-1.367,0.206) (-1.372,0.195) (-1.377,0.185) (-1.381,0.174) (-1.385,0.164) (-1.389,0.153) (-1.392,0.143) (-1.395,0.132) (-1.398,0.121) (-1.401,0.110) (-1.404,0.099) (-1.406,0.088) (-1.408,0.077) (-1.410,0.066) (-1.411,0.055) (-1.412,0.044) (-1.413,0.033) (-1.414,0.022) (-1.414,0.011) (-1.414,-0.000) (-1.414,-0.011) (-1.414,-0.022) (-1.413,-0.033) (-1.412,-0.044) (-1.411,-0.055) (-1.410,-0.066) (-1.408,-0.077) (-1.406,-0.088) (-1.404,-0.099) (-1.401,-0.110) (-1.398,-0.121) (-1.395,-0.132) (-1.392,-0.143) (-1.389,-0.153) (-1.385,-0.164) (-1.381,-0.174) (-1.377,-0.185) (-1.372,-0.195) (-1.367,-0.206) (-1.362,-0.216) (-1.357,-0.226) (-1.351,-0.236) (-1.346,-0.246) (-1.340,-0.256) (-1.333,-0.265) (-1.327,-0.275) (-1.320,-0.284) (-1.313,-0.293) (-1.306,-0.303) (-1.298,-0.312) (-1.290,-0.321) (-1.282,-0.329) (-1.274,-0.338) (-1.266,-0.346) (-1.257,-0.354) (-1.248,-0.363) (-1.239,-0.370) (-1.229,-0.378) (-1.220,-0.386) (-1.210,-0.393) (-1.200,-0.400) (-1.189,-0.407) (-1.179,-0.414) (-1.168,-0.421) (-1.157,-0.427) (-1.146,-0.433) (-1.134,-0.439) (-1.123,-0.444) (-1.111,-0.450) (-1.099,-0.455) (-1.086,-0.460) (-1.074,-0.465) (-1.061,-0.469) (-1.048,-0.473) (-1.035,-0.477) (-1.022,-0.481) (-1.008,-0.484) (-0.994,-0.487) (-0.980,-0.490) (-0.966,-0.492) (-0.952,-0.494) (-0.937,-0.496) (-0.922,-0.498) (-0.907,-0.499) (-0.892,-0.499) (-0.876,-0.500) (-0.861,-0.500) (-0.845,-0.500) (-0.829,-0.499) (-0.812,-0.498) (-0.796,-0.496) (-0.779,-0.494) (-0.762,-0.492) (-0.745,-0.489) (-0.727,-0.486) (-0.710,-0.482) (-0.692,-0.478) (-0.673,-0.473) (-0.655,-0.468) (-0.636,-0.462) (-0.617,-0.456) (-0.597,-0.449) (-0.577,-0.441) (-0.557,-0.432) (-0.537,-0.423) (-0.515,-0.413) (-0.494,-0.402) (-0.472,-0.390) (-0.449,-0.377) (-0.425,-0.363) (-0.401,-0.348) (-0.376,-0.331) (-0.349,-0.313) (-0.321,-0.292) (-0.291,-0.269) (-0.258,-0.243) (-0.222,-0.212) (-0.180,-0.174) (-0.126,-0.124) (-0.000,-0.000)};
  % points reguliers : tangente (orange) et gradient unitaire x 0.45 (vert)
  \draw[couleurD,thick] (1.4142,-0.3) -- (1.4142,0.3);
  \draw[couleurZ,thick,->,>=Stealth] (1.4142,0) -- (1.8642,0);
  \draw[couleurD,thick] (0.566,0.5) -- (1.166,0.5);
  \draw[couleurZ,thick,->,>=Stealth] (0.866,0.5) -- (0.866,0.95);
  \draw[couleurD,thick] (0.8216,-0.5698) -- (1.3758,-0.3402);
  \draw[couleurZ,thick,->,>=Stealth] (1.0987,-0.4551) -- (1.2709,-0.8709);
  \draw[couleurD,thick] (-1.166,-0.5) -- (-0.566,-0.5);
  \draw[couleurZ,thick,->,>=Stealth] (-0.866,-0.5) -- (-0.866,-0.95);
  \foreach \p in {(1.4142,0),(0.866,0.5),(1.0987,-0.4551),(-0.866,-0.5)} \fill[couleurN] \p circle (0.8pt);
  \fill[rougecours] (0,0) circle (1.0pt);
  \node[rougecours,font=\small,anchor=north,fill=white,inner sep=1pt] at (0,-0.9) {$\grad g(0, 0) = 0$};
  \node[couleurN,font=\small,anchor=east] at (-0.95,0.62) {$(x^{2} + y^{2})^{2} = 2(x^{2} - y^{2})$};
\end{tikzpicture}
\end{center}

\noindent La lemniscate de Bernoulli, en bleu, a pour équation $g(x, y) = 0$, avec
\[ g(x, y) = (x^{2} + y^{2})^{2} - 2(x^{2} - y^{2}), \qquad \grad g(x, y) = 4\big(x(x^{2} + y^{2} - 1), y(x^{2} + y^{2} + 1)\big). \]
En chaque point régulier, le gradient, en vert, est perpendiculaire à la tangente, en orange : au point $(\sqrt{2}, 0)$, la tangente est verticale ; au point $\left(\dfrac{\sqrt{3}}{2}, \dfrac{1}{2}\right)$, où l'ordonnée est maximale, elle est horizontale. À l'origine, le gradient est nul, et la courbe passe deux fois, avec les deux tangentes $y = \pm x$, en pointillés rouges : l'ensemble des vecteurs tangents contient ces deux droites, et ce n'est donc pas une droite.

\subsection{Exemples}

\begin{exemples}
  \item \emph{La sphère.} Pour $g(x) = \norme{x}_{2}^{2} - R^{2}$, avec $R > 0$, on a $\grad g(a) = 2a \neq 0$ sur la sphère : tous ses points sont réguliers, et l'hyperplan tangent en $a$ a pour équation $\ps{a}{x - a} = 0$, soit $\ps{a}{x} = R^{2}$. Il est orthogonal au rayon, comme en géométrie élémentaire.
  \item \emph{Un graphe.} Soit $\Omega$ un ouvert de $\R^{2}$, $h : \Omega \to \R$ de classe $\Cc^{1}$, et $S$ son graphe, d'équation $g(x, y, z) = h(x, y) - z = 0$. Comme $\grad g = \left(\dpd{h}{x}, \dpd{h}{y}, -1\right)$ ne s'annule jamais, tous les points sont réguliers, et le plan tangent en $(a, h(a))$ est
  \[ z = h(a) + \dpd{h}{x}(a)\,(x - a_{1}) + \dpd{h}{y}(a)\,(y - a_{2}). \]
  On retrouve le plan tangent du chapitre 2 : pour $h(x, y) = 2{,}5 - \dfrac{x^{2} + 2y^{2}}{4}$ et $a = \left(1, -\dfrac{1}{2}\right)$, on obtient $z = \dfrac{17}{8} - \dfrac{x - 1}{2} + \dfrac{1}{2}\left(y + \dfrac{1}{2}\right)$.
  \item \emph{L'hyperboloïde à une nappe.} Pour $g(x, y, z) = x^{2} + y^{2} - z^{2} - 1$, on a $\grad g = 2(x, y, -z)$, qui ne s'annule pas sur $S$, puisque $(0, 0, 0) \notin S$. Au point $a = (1, 1, 1)$, $\grad g(a) = (2, 2, -2)$ et le plan tangent a pour équation $x + y - z = 1$.
\end{exemples}

% ---- hyperboloide x^2 + y^2 - z^2 = 1 ; a = (1,1,1) ; plan tangent x + y - z = 1 ;
% ---- droites de S passant par a : a + t(1,0,1) et a + t(0,1,1) (g s'annule identiquement, verifie avec sympy)
\begin{center}
\begin{tikzpicture}
  \begin{axis}[width=12.5cm,height=11cm,view={135}{14},axis lines=none,axis equal image,clip=false,
    z buffer=sort,colormap={bleu}{color=(couleurN!25) color=(couleurN!70)}]
    \addplot3[surf,shader=faceted interp,draw=couleurN!60,line width=0.1pt,opacity=0.75,samples=26,samples y=11,domain=0:360,y domain=-1.6:1.6]
      ({sqrt(1+y^2)*cos(x)},{sqrt(1+y^2)*sin(x)},{y});
    \addplot3[surf,shader=flat,fill=couleurD!35,draw=couleurD!60,opacity=0.55,samples=2,samples y=2,domain=-0.9:0.9,y domain=-0.9:0.9]
      ({1+x},{1+y},{1+x+y});
    \addplot3[couleurN!80!black,very thick,domain=-1.4:0.6,samples=2] ({1+x},{1},{1+x});
    \addplot3[couleurN!80!black,very thick,domain=-1.4:0.6,samples=2] ({1},{1+x},{1+x});
    \addplot3[couleurZ,very thick,-{Stealth}] coordinates {(1,1,1) (1.55,1.55,0.45)};
    \addplot3[only marks,mark=*,mark size=1.6pt,color=rougecours] coordinates {(1,1,1)};
    \node[rougecours,font=\small,anchor=south east] at (axis cs:1,1,1.05) {$a$};
    \node[couleurZ,font=\small,anchor=north west] at (axis cs:1.55,1.55,0.45) {$\grad g(a)$};
    \node[couleurD!80!black,font=\small,anchor=west,fill=white,inner sep=1pt] at (axis cs:0.1,1.9,1.0) {$x + y - z = 1$};
  \end{axis}
\end{tikzpicture}
\end{center}

\noindent L'hyperboloïde, en bleu, et son plan tangent en $a = (1, 1, 1)$, en orange, ont pour vecteur normal commun le gradient, en vert. Le plan coupe la surface selon deux droites, en bleu foncé, dirigées par $(1, 0, 1)$ et $(0, 1, 1)$ : on vérifie que $g(1 + t, 1, 1 + t) = g(1, 1 + t, 1 + t) = 0$ pour tout $t$. Un plan tangent peut donc rencontrer la surface bien au-delà du point de contact ; il ne la «~touche~» qu'au premier ordre.

\subsection{Intersections d'hypersurfaces}

Plusieurs équations définissent une partie plus petite : dans $\R^{3}$, deux surfaces se coupent en général selon une courbe. Le cas de la sphère et du plan, rencontré à la section 3.4, se généralise.

\begin{cadre}
\propo{11} Soit $g = (g_{1}, \dots, g_{n}) : U \to \R^{n}$ de classe $\Cc^{1}$, avec $n < p$, $S = \{x \in U \mid g(x) = 0\}$, et $a \in S$ tel que les gradients $\grad g_{1}(a)$, \dots, $\grad g_{n}(a)$ soient linéairement indépendants. Alors, au voisinage de $a$, $S$ est le graphe d'une fonction de classe $\Cc^{1}$ de $p - n$ des coordonnées, et
\[ \rmath{T_{a}S = \bigcap_{i = 1}^{n} \grad g_{i}(a)^{\perp} = \Ker \dd g(a)}, \]
sous-espace de dimension $p - n$.
\end{cadre}

\begin{demo}
La jacobienne $\Jac_{g}(a)$, dont les lignes sont les gradients, est de rang $n$ : elle possède $n$ colonnes linéairement indépendantes. Quitte à renuméroter les coordonnées, ce sont les $n$ dernières, et l'on écrit $x = (x', x'')$ avec $x' \in \R^{p - n}$ et $x'' \in \R^{n}$. Le bloc $\Jac_{x''} g(a)$ est inversible, et la proposition 8 donne $\varphi$ de classe $\Cc^{1}$ telle que $S$ coïncide près de $a$ avec le graphe $\{(x', \varphi(x'))\}$.

Si $\gamma$ est une courbe tracée sur $S$, dérivable en $0$, avec $\gamma(0) = a$, alors $g \circ \gamma = 0$, d'où $\dd g(a)\cdot\gamma'(0) = 0$ : $T_{a}S \subset \Ker \dd g(a)$. Réciproquement, soit $v = (v', v'') \in \Ker \dd g(a)$, c'est-à-dire $\Jac_{x'} g(a)\,v' + \Jac_{x''} g(a)\,v'' = 0$. La proposition 8 donne $\Jac_{\varphi}(a')\,v' = -\Jac_{x''} g(a)^{-1}\Jac_{x'} g(a)\,v' = v''$, et la courbe $\gamma(t) = (a' + tv', \varphi(a' + tv'))$, tracée sur $S$, a pour vecteur dérivé $v$ en $0$. Enfin, $\dd g(a)$ est surjective, puisque son rang vaut $n$, et le théorème du rang donne $\dim \Ker \dd g(a) = p - n$.
\end{demo}

\begin{exemple}
Pour la courbe $x^{2} + y^{2} + z^{2} = 3$, $x + y - z = 1$ de la section 3.4, les gradients en $(1, 1, 1)$ sont $(2, 2, 2)$ et $(1, 1, -1)$, indépendants. La tangente est dirigée par un vecteur orthogonal aux deux, par exemple leur produit vectoriel, $(2, 2, 2) \wedge (1, 1, -1) = (-4, 4, 0)$, colinéaire à $(1, -1, 0)$ : c'est le vecteur trouvé par le calcul de $y'(1)$ et $z'(1)$.
\end{exemple}

\begin{remarque}
\emph{Vers le chapitre 5.} Si $f : U \to \R$ de classe $\Cc^{1}$ atteint en un point régulier $a$ de $S$ son maximum sur $S$, alors, pour toute courbe $\gamma$ tracée sur $S$ avec $\gamma(0) = a$, la fonction $f \circ \gamma$ a un maximum en $0$, donc $\dd f(a)\cdot\gamma'(0) = 0$ : le gradient $\grad f(a)$ est orthogonal à $T_{a}S$, donc colinéaire à $\grad g(a)$. C'est le principe des multiplicateurs de Lagrange, que le chapitre 5 développera.
\end{remarque}
\newpage
%===============================================================================
\section{Bilan}
%===============================================================================

Le fil du chapitre est la résolution d'équations par linéarisation. La méthode de Newton remplace l'équation par l'équation affine tangente ; en figeant la différentielle, on obtient une contraction, et le théorème du point fixe du chapitre 1 démontre que l'équation $f(x) = y$ a une unique solution près de $a$ dès que $\dd f(a)$ est inversible et que $f$ est de classe $\Cc^{1}$. C'est le théorème d'inversion locale, dont le théorème des fonctions implicites est une reformulation, et qui décrit, près d'un point régulier, une hypersurface comme un graphe ayant un hyperplan tangent.

\begin{center}
\renewcommand{\arraystretch}{1.35}
\begin{tabularx}{15cm}{|>{\raggedright\arraybackslash}p{4.3cm}|>{\raggedright\arraybackslash}X|>{\centering\arraybackslash}p{2.4cm}|}
\hline
\rowcolor{couleurD!35} \textbf{Situation} & \textbf{Résultat} & \textbf{Référence} \\ \hline
$x_{0} > 0$, $a > 0$ & Héron : $x_{n + 1} - \sqrt{a} = \dfrac{(x_{n} - \sqrt{a})^{2}}{2x_{n}}$ & proposition 1 \\ \hline
$E$, $F$ de Banach & $\Isom(E, F)$ ouvert, inversion continue & proposition 2 \\ \hline
$f(c) = 0$, $\dd f(c)$ inversible, $\dd f$ lipschitzienne & Newton : $\norme{x_{n + 1} - c} \leqslant C\norme{x_{n} - c}^{2}$ & proposition 3 \\ \hline
$\varphi$ $k$-lipschitzienne, $k < 1$ & $\id + \varphi$ injective, image contenant $\Bf(0, (1 - k)r)$ & proposition 4 \\ \hline
$\nop{L \circ (\dd f(x) - \dd f(a))} \leqslant \dfrac{1}{2}$ sur $\Bf(a, r)$ & jacobienne figée : solution unique, erreur $\leqslant 2^{-n}r$ & proposition 5 \\ \hline
$f$ difféomorphisme & $\dd(f^{-1})(f(x)) = \dd f(x)^{-1}$ & proposition 6 \\ \hline
$f$ de classe $\Cc^{1}$, $\dd f(a)$ inversible & difféomorphisme local & théorème 1 \\ \hline
de plus $f$ injective, $\dd f$ partout inversible & difféomorphisme global sur $f(U)$ & proposition 7 \\ \hline
$f(a, b) = 0$, $\dd_{2} f(a, b)$ inversible & $y = \varphi(x)$ près de $(a, b)$, $\dd\varphi = -\dd_{2} f^{-1} \circ \dd_{1} f$ & théorème 2, proposition 8 \\ \hline
$g(a) = 0$, $\grad g(a) \neq 0$ & $S$ localement un graphe, $T_{a}S = \grad g(a)^{\perp}$ & proposition 9, théorème 3 \\ \hline
$\grad f(a) \neq 0$ & gradient orthogonal à la ligne de niveau & proposition 10 \\ \hline
gradients $\grad g_{i}(a)$ indépendants & $T_{a}S = \Ker \dd g(a)$, de dimension $p - n$ & proposition 11 \\ \hline
\end{tabularx}
\end{center}

\begin{center}
\begin{tikzpicture}[q/.style={draw=couleurN,fill=couleurN!8,rounded corners=3pt,align=center,font=\small,minimum height=1.05cm,text width=2.3cm},
  fl/.style={-{Stealth},thick,gray}]
  \node[q] (P) at (0,0) {point fixe\\ (chapitre 1)};
  \node[q] (N) at (3.05,0) {Newton à\\ jacobienne figée};
  \node[q] (I) at (6.1,0) {inversion\\ locale};
  \node[q] (F) at (9.15,0) {fonctions\\ implicites};
  \node[q] (H) at (12.2,0) {hyperplan\\ tangent};
  \node[q] (A) at (3.05,-1.9) {accroissements\\ finis\\ (chapitre 2)};
  \node[q] (G) at (6.1,-1.9) {$\Isom(E, F)$\\ ouvert};
  \node[q] (L) at (12.2,-1.9) {Lagrange\\ (chapitre 5)};
  \draw[fl] (P) -- (N);
  \draw[fl] (N) -- (I);
  \draw[fl] (I) -- (F);
  \draw[fl] (F) -- (H);
  \draw[fl] (A) -- (N);
  \draw[fl] (G) -- (I);
  \draw[fl,densely dashed] (H) -- (L);
\end{tikzpicture}
\end{center}

\noindent Chaque flèche pleine indique un résultat qui sert à démontrer le suivant ; la flèche en tirets annonce le chapitre 5. La complétude intervient à un seul endroit, le point fixe, et la classe $\Cc^{1}$ à un seul endroit aussi, le choix d'une boule où la différentielle reste proche de $\dd f(a)$.

La suite du cours approfondit ces résultats. Le chapitre 4 définit les différentielles d'ordre supérieur, qui mesurent l'écart entre une fonction et son approximation affine, et donne les formules de Taylor ; elles expliqueront la convergence quadratique de la méthode de Newton pour les fonctions de classe $\Cc^{2}$. Le chapitre 5 cherche les extrema, libres puis liés : sur une hypersurface, le théorème 3 conduit aux multiplicateurs de Lagrange. Le chapitre 6, enfin, réutilise le point fixe pour résoudre les équations différentielles.

\end{document}
